\documentclass[ecta,nameyear,draft]{econsocart}
\usepackage{amsmath,amsfonts,amssymb,amsthm,booktabs,array,longtable}
\usepackage[T1]{fontenc}
\usepackage{mathpazo,microtype,needspace,xurl}
\usepackage{xr-hyper}
\RequirePackage[colorlinks,citecolor=blue,linkcolor=blue,urlcolor=blue]{hyperref}
\startlocaldefs
\newtheorem{theorem}{Theorem}
\newtheorem{proposition}[theorem]{Proposition}
\newtheorem{lemma}[theorem]{Lemma}
\newtheorem{corollary}[theorem]{Corollary}
\newcommand{\E}{\mathbb E}
\newcommand{\R}{\mathbb R}
\endlocaldefs
\setlength{\emergencystretch}{3em}
\makeatletter\def\copyright@text{Prepared for referee review}\makeatother

\setcounter{secnumdepth}{2}

% Stable current-document identity for clean typesetting reproduction.
\ifdefined\pdfinfoomitdate\pdfinfoomitdate=1\fi
\ifdefined\pdfsuppressptexinfo\pdfsuppressptexinfo=15\fi
\ifdefined\pdftrailerid\pdftrailerid{}\fi

\externaldocument[][nocite]{build/supp}[supp.pdf]
\externaldocument[][nocite]{build/appendix}[appendix.pdf]
\begin{document}
\begin{frontmatter}
\title{Neural Bellman Operators}
\runtitle{Neural Bellman Operators}
\begin{aug}
\author[id=au1,addressref={add1}]{\fnms{Qian}~\snm{QI}\ead[label=e92ECTA]{qiqian@pku.edu.cn}}
\address[id=add1]{Peking University}
\end{aug}
\begin{abstract}
This paper develops Neural Bellman Operators for constructing feasible dynamic policies and qualifying their loss against the optimum under the original economic law. Own-future continuation construction retains feasible action witnesses and an explicit approximation budget. Shared-readout and occupation restrictions extend the account to complete policies. Under recursive risk, centering the conditional dual at the stored continuation yields state-dependent bounds and safe kernel exclusions. A continuous investment application derives primitive curvature, global value supports and all-action certificates without replacing the innovation law by a discrete model. Six frozen continuation fits are compared with conventional online risk dynamic programming at a common original-optimum tolerance. ReLU fits have narrower whole-state residual brackets, while all policy-cost contrasts identify equivalence at the prespecified margin and none identifies strict superiority. Stored continuations reduce measured evaluation work in this experiment; paid screening does not. The joint-domain construction, horizon-resource account and broader economic applications retain their own hypotheses. These results connect a constructed approximation to its implemented economic decision without equating certificate sharpness with algorithmic superiority.
\end{abstract}
\begin{keyword}
\kwd{Dynamic programming}\kwd{Neural approximation}\kwd{Policy certification}
\kwd{Information acquisition}\kwd{Computational economics}
\end{keyword}
\end{frontmatter}
\section{Introduction}\label{sec:introduction65}
A continuation approximation matters through the economic decision it produces. A small prediction error, a feasible controller, a bound on original-optimum policy loss and a saving in computation are different achievements. This paper develops Neural Bellman Operators (NBOs) to connect them. The benchmark is the original economic law and objective, and the object evaluated is the complete implemented policy, including its observation rule and numerical action realization.

An NBO fits each date against its own attained future. Our constructive backend uses primitive state and action covers, certified integration and feasible action witnesses to return an affine--ReLU continuation and an explicit original-optimum error budget. It does not assume that unrestricted nonconvex fitting succeeds. On a tensor cover, original-owner compilation preserves both the represented function and the owner's feasible action at off-grid states. A complementary recursive implementation uses verified curvature and protected conventional refinement, allowing trained proposals to change optional query locations without weakening the economic certificate.

The central comparison has the form
\[
 \underline V\ \leq\ V^*\ \leq\ J^{\widehat\pi}\ \leq\ \overline J.
\]
Its lower side must cover every original feasible action. Its upper side must concern the actual controller and original preference, not only its fitted objective. We first obtain these inequalities from own-future residuals and explicit acquisition error. Shared coefficients then restrict which acquired actions can occur jointly, and complete-policy occupations propagate changed decisions through their own state distributions. The corresponding finite actionwise, compatibility and exact tolerance results are retained in the numerical appendix rather than reproduced as a sequence of separate main-text studies.

Under nonlinear recursive preferences, a physical-expectation identity is generally not available. We use monotonicity and cash invariance for the policy comparison. Centering the original conditional-risk dual at the stored continuation preserves the nonnegative cost of changing the risk kernel. State-dependent error functions can then yield a substantially smaller loss bound than a scalar worst-state charge. The exact finite risk-occupation representation and strict support exclusions implement this idea under their stated hypotheses. A frozen reference maximizer is not substituted for the original nonlinear risk evaluation, and equality at an exclusion boundary is retained.

A general comparison theorem is not yet an executable numerical method. The continuous investment application supplies the missing primitive construction. The original transition is nonlinear in the state, the feasible capacity depends on that state, actions range over an interval and shocks have a continuous law. We prove joint curvature for the one-date risk objective, derive global affine value supports with an explicit multiplier-stationarity allowance, and establish a state modulus. These results yield an all-action lower witness and a complete-policy upper cost under the original continuous innovation law. Innovation bins are sets on which inequalities hold, not new point masses in the economic model.

This application also permits construction, learning, qualification, computation and policy-cost inference to be evaluated on the same problem. At each of two fixed risk levels, oblique ReLU, additive spline and projected-quadratic continuations use the same labels and nominal feature counts; conventional online risk dynamic programming pays for its original continuation queries. Both an own-continuation target and a common original-optimum target are required. Training and calibration are separated from the 4,096 initial-state observations used in the evaluation, and all 16,384 arm-state records are retained. The ReLU fits have narrower whole-square residual brackets, but every simultaneous policy-cost interval includes zero. All four policies are equivalent at the prespecified normalized margin; no strict cost superiority is identified.

Stored continuations reduce measured work on this workload even after paying for the original-value lower witness. This is not a uniquely neural result: the conventional fitted representations obtain comparable reductions, and the projected-quadratic features have lower effective rank. Likewise, strict continuum screening removes candidate bins but is slightly slower once its proof is paid for. The article distinguishes these observations from analytical certificate improvements, from construction overhead and from a full cross-architecture or high-dimensional implementation claim.

\paragraph{Relation to existing methods.}
Lipschitz extensions and distance transforms provide established construction tools \citep{mcshane1934,felzenszwalb2012}. Approximate policy iteration, safe improvement and convex dynamic programming have established error and implementation analyses \citep{munos2003,pirotta2013,scherrer2015,yang2020convex}. Learned warm starts and optimization-based control address related computational questions \citep{sakaue2023warm,sambharya2023,agrawal2020control}. Our construction retains the continuation's own future, the feasible action owner and the original-optimum accounting obligations. An exact neural encoding and a native encoding of the same function are not two different economic methods.

Risk-averse Bellman equations and policy iteration are established results \citep{ruszczynski2010risk,ruszczynski2014erratum}; feature-based risk evaluation and structured improvement are also established research directions \citep{ruszczynskiyang2025}. We do not claim conditional-risk duality or elementary centering as new. The theorem-level comparison in Section~\ref{sec:comparison95} distinguishes those tools from the state-dependent qualification object, its risk-distorted occupation representation and its constructive all-action implementation. Standard convex duality supplies the support principle \citep{boyd2004}; independently verified multipliers, the original feasible menu and the paid error set determine whether it licenses an exclusion here.

Recent feature-based risk learning also includes multipattern risk approximation and high-probability regret analysis \citep{ruszczynskizhang2026}. Those results concern learning in finite-state risk models under their feature and sampling assumptions. Our fixed-accuracy horizon construction instead uses the specified continuous investment primitives and pays for deterministic state, action and innovation resolution. A rational-operation bound for that construction is not an episode-regret bound and does not supersede the learning result.

The retained economic accuracy and selection results similarly use established concentration, control-variate and decision-theoretic tools \citep{munos2006variance,jiang2016,howard2020time,karampatziakis2021,manski2004,floyd1962}. Their role is to connect verified optimization accuracy to a sufficient economic stopping region. A narrower certificate for an unchanged controller is not a new economic observation. The full earlier investment catalogue, including adverse rankings, unresolved comparisons and the effect of affine calibration in both ReLU and Bernstein bases, remains in the appendix.

\paragraph{Scope and reading path.}
Sections~\ref{sec:core49} and \ref{sec:construction49} state the economic operator and feasible own-future construction. Section~\ref{sec:accuracy62} develops primitive original-optimum accuracy. Complete-policy, shared-domain and nonlinear-risk results follow in Sections~\ref{sec:recursive86}, \ref{sec:population86}, \ref{sec:risk91} and \ref{sec:centeredrisk92}. Sections~\ref{sec:continuous95} and \ref{sec:study95} give the continuous primitive construction and its frozen investment comparison. The stored-successor resource and prospective-decision results retain explicit acquisition, evaluation and externally supplied economic charges.

The Technical Supplement contains complete proofs. The Numerical and Development Appendix preserves the specialized qualification results, all earlier empirical catalogues, singular-transition arguments, recursive-query and vector-control evidence, and the complete new numerical comparisons. Preserved companion editions retain diffusions, recursive and endogenous preferences, temporal selves and games under their original hypotheses. Expected-cost identities are not transferred to those objectives without their own comparison principles. The present numerical application strengthens the original NBO program; it does not replace that program by a finite-menu or qualification-only subject.

\section{The economic operator and its policy account}\label{sec:core49}
Consider dates $t=0,\ldots,T$ and an invariant compact state domain $X\subset\R^d$. At date $t$, the action belongs to a nonempty compact set $A_t(x)$. The transition is $F_t(x,a,Z_{t+1})$, the current cost is $c_t(x,a)$, and the terminal cost is $g(x)$. Write $B_0=1$ and $B_t=\prod_{j<t}\beta_j$, with $\beta_j>0$. An implemented policy includes its observation rule, action selection and feasible numerical realization, not just its stored parameters.

For a continuation $v$, define
\begin{equation}\label{eq:operator49}
 Q_t(v;x,a)=c_t(x,a)+\beta_t\rho_t(v(F_t(x,a,Z))),\qquad
 \mathcal T_tv(x)=\inf_{a\in A_t(x)}Q_t(v;x,a).
\end{equation}
Throughout the deterministic comparison, the evaluations $\rho_t$ are monotone and cash invariant on a common domain containing all functions and constant translates used below. They are therefore nonexpansive in the supremum norm. Continuous objectives, effective compact feasible correspondences and measurable minimizers, or measurable arbitrarily near-minimizers with their error charged, are assumed. The stipulated conditional recursion gives Bellman comparison against all admissible adapted controls. In the earlier expected-cost experiments $\rho_t$ is conditional expectation, costs are bounded and innovations have their original continuous uniform law. Section~\ref{sec:continuous95} evaluates the same investment technology under an explicitly nested AVaR preference and retains the original continuous innovation law. The direct-cost inference in Section~\ref{sec:decisions49} is specific to expectation; it is not automatically an inference theorem for nonlinear recursive utility or games.

A Neural Bellman Operator is a construction and evaluation procedure for \eqref{eq:operator49}. It forms targets using its \emph{own} fitted $f_{t+1}$, returns a neural continuation and a feasible policy, and records an economic error allowance. It does not receive the optimal continuation as a training label. Arbitrary trained networks can be evaluated by the first result below. The subsequent constructive backend specifies how a network and its actor are obtained without assuming successful nonconvex optimization.

\subsection{Action-dependent errors and complete-policy loss}
Let $H(a)$ be a true current cost and $\widehat H(a)$ its continuation approximation. If $\widehat a$ is $\delta$-optimal for $\widehat H$ on the same feasible action set, then
\begin{equation}\label{eq:center49}
 H(\widehat a)-\inf_aH(a)\le\delta+
 \operatorname{osc}_{a}\{\widehat H(a)-H(a)\}.
\end{equation}
Indeed, compare $\widehat a$ with an arbitrarily near-optimal true action and cancel the approximate objectives. The infimum limit gives the result. Additive continuation levels cancel; action-dependent discrepancies do not. Under a change in future policies, the continuation must be transported to the new economic problem before this inequality licenses reuse. The full development edition retains the original quantitative transport, recursive-domain and economic-application results. No transport bound is presumed valid after an unaccounted change of constraints or preferences.

\begin{proposition}[One-sided policy sandwich]\label{prop:sandwich49}
Suppose, uniformly in state,
\[
 f_t-\mathcal T_tf_{t+1}\le u_t,\qquad
 \mathcal T_t^\pi f_{t+1}-f_t\le v_t,
 \qquad \ell_T\le f_T-g\le u_T.
\]
Then the actual feasible policy satisfies
\begin{equation}\label{eq:sandwich49}
 0\le J_0^\pi-V_0^*\le
 \sum_{t<T}B_t(u_t+v_t)+B_T(u_T-\ell_T).
\end{equation}
The bound is invariant to date-specific constant shifts of the fitted continuations, provided the associated residual accounts are shifted consistently.
\end{proposition}
\noindent The proof is in Supplement, Section~\ref{supp:proof-core49-1}.


A sharper valid residual account for the same continuation and the same deployed actor sharpens knowledge of the policy. It does not retrain that policy or change its realized cost. Conversely, subtracting two upper bounds in \eqref{eq:sandwich49} does not identify the difference in actual costs. These distinctions govern both the historical comparisons and the new experiments.

\section{Constructing and compiling a feasible neural policy}\label{sec:construction49}
Assume the primitive objective has, for feasible state--action pairs, state modulus $A_t^Q=C_{x,t}+\beta_tF_{x,t}L_{t+1}$ and action modulus $D_t=C_{a,t}+\beta_tF_{a,t}L_{t+1}$. Let
\[
 d_H(A_t(x),A_t(z))\le K_t^A\|x-z\|_1,
 \qquad L_t=A_t^Q+D_tK_t^A,\qquad L_T=L_g.
\]
The Hausdorff condition and the primitive moduli imply that $\mathcal T_tf_{t+1}$ is $L_t$-Lipschitz whenever $f_{t+1}$ is $L_{t+1}$-Lipschitz. All moduli and numerical integration procedures are supplied by the economic primitives, not fitted from successful outcomes.

At state nodes $x_{t,i}$ with covering radius $h_t$, form effective feasible $k_t$-nets. Evaluate their own-future objectives with certified error at most $e_t$. Minimize the computed labels and retain the corresponding \emph{original} action $a_{t,i}$. A Borel repair $\Gamma_{t,i}(x)$ is feasible at $x$ and satisfies
\begin{equation}\label{eq:repair49}
 \|\Gamma_{t,i}(x)-a_{t,i}\|_1\le K_t^A\|x-x_{t,i}\|_1+\zeta_t.
\end{equation}
The allowance $\zeta_t$ concerns displacement; feasibility itself may not fail. Define
\begin{equation}\label{eq:cone49}
 f_t(x)=\min_i\{y_{t,i}+L_t\|x-x_{t,i}\|_1\},\quad
 i_t(x)=\min\arg\min_i\{y_{t,i}+L_t\|x-x_{t,i}\|_1\}.
\end{equation}
Return $\pi_t(x)=\Gamma_{t,i_t(x)}(x)$. Terminal labels approximate $g$ with error $e_T$ and use the same cone construction with modulus $L_g$.

\begin{theorem}[Own-future construction with feasible witnesses]\label{thm:construct49}
The functions in \eqref{eq:cone49} have exact affine--ReLU realizations and Lipschitz moduli $L_t$ independent of label error. The returned policy is feasible and measurable. It satisfies
\begin{align}\label{eq:enclosures49}
 -e_t\le f_t-\mathcal T_tf_{t+1}&\le D_tk_t+e_t+2L_th_t,\\
 \mathcal T_t^\pi f_{t+1}-f_t&\le e_t+D_t\zeta_t,\notag
\end{align}
and hence
\begin{equation}\label{eq:constructionloss49}
 J_0^\pi-V_0^*\le
 \sum_{t<T}B_t(D_tk_t+2e_t+2L_th_t+D_t\zeta_t)
 +B_T(2e_T+2L_gh_T).
\end{equation}
With effective covers, repairs and certified queries at the prescribed accuracy, allocating each nonnegative term a share of a positive economic tolerance gives a terminating finite construction. A fixed numerical precision or a capped finite experimental ladder need not contain that allocation.
\end{theorem}
\noindent The proof is in Supplement, Section~\ref{supp:proof-core49-2}.


The Lipschitz-extension principle is classical \citep{mcshane1934}. A native min-plus envelope and its exact ReLU circuit, using the same labels, witnesses and ties, are \emph{one method with two encodings}. They return the same function and policy. Neither Theorem~\ref{thm:construct49} nor a faster representation proves a uniquely neural capability or a training-success probability for Adam or L-BFGS. The mathematical increment is the complete feasible, own-future construction and implementation account; the empirical question is its work and policy value relative to conventional generators.

\subsection{Original-owner compilation on a tensor cover}
On a tensor cover $\mathcal X=\prod_{j=1}^d\{x_{j,0},\ldots,x_{j,n_j}\}$, put $S=\prod_j(n_j+1)$. For arbitrary labels $y_i$, including inconsistent inexact labels, compute at each grid node $v$
\[
 (z_v,o_v)=\min_{\rm lex,i}(y_i+L\|v-x_i\|_1,i).
\]
The index $o_v$ is the original label owner, not the transformed node $v$.

\begin{theorem}[Continuous original-owner identity]\label{thm:compile49}
For a closed tensor cell $C$ with corners $\mathcal V(C)$, and $x\in C$,
\begin{align}\label{eq:compile49}
 f(x)&=\min_{v\in\mathcal V(C)}(z_v+L\|x-v\|_1),\\
 (f(x),i_*(x))&=\min_{\rm lex,v\in\mathcal V(C)}
 (y_{o_v}+L\|x-x_{o_v}\|_1,o_v).\notag
\end{align}
Coordinatewise forward and backward relaxations compute $(z,o)$ in $O(dS)$ additions and comparisons and $O(S)$ storage. After locating the cell, a query uses $O(d2^d)$ arithmetic operations. Exact rational comparisons preserve all ties, including $L=0$. Using the original owner's repaired action preserves the uncompiled policy, its cost and its certificate exactly.
\end{theorem}
\noindent The proof is in Supplement, Section~\ref{supp:proof-core49-3}.


The nodal algorithm is the classical distance transform \citep{felzenszwalb2012}; preserving the original action through continuous off-grid queries is the policy obligation. If there are $J_{t,i}$ node actions and $M_t$ integration representatives, write $Q_t=M_t\sum_iJ_{t,i}$. Literal dense cone evaluation costs $O(\sum_t dS_{t+1}Q_t)$ in representation work. Compilation replaces this by $O(\sum_t[dS_{t+1}+Q_td2^d])$, plus cell location, primitive evaluation, integration, feasible repair, certificate and output. Rational bit lengths and state boxes crossing many cells incur additional costs. There is no dimension-free complexity assertion.

\subsection{Acquired states and conventional actors}
A $b$-bit coordinate sensor reports a dyadic cell and its midpoint $\widehat x$. On the unit cube, $\|x-\widehat x\|_1\le r(b)=d2^{-b-1}$. For the investment capacity $\kappa(x)$ below, clip the original action to the smallest capacity throughout that cell and round down to an action quantum $\Delta_t$. This is feasible at every state in the cell. Relative to an exact true-state repair, its displacement is at most $2K_t^Ar(b)+\Delta_t$. Transporting the witness cone from $\widehat x$ to $x$ adds at most $2L_tr(b)$, plus any certified selector objective error $\nu_t$. Thus the implemented-policy account adds
\begin{equation}\label{eq:acquire49}
 \sum_{t<T}B_t\{(2L_t+2D_tK_t^A)r(b)+D_t\Delta_t+\nu_t\}.
\end{equation}
No continuity of the discrete owner selector is assumed. At uncertain numerical comparisons the simulator retains all possible actors, rather than choosing a favorable one.

For conventional FVI the continuation is multilinear on its realized tensor grid. Its actual modulus is bounded by the largest coordinate edge slope. A nearest-node actor uses a feasible node minimizer and true-state repair. Let $h_t$ be the largest half-cell $\ell^1$ radius, $L_t^{\rm gr}$ the Bellman graph modulus, and $e_t$ the node query allowance. Its optimal residual upper bound is $D_tk_t+e_t+L_t^{\rm gr}h_t$. Its selected-policy upper bound is $e_t+E_t$, where
\[
 E_t=\sup_x\{y_{i(x)}+L_t^{\rm gr}\|x-x_{i(x)}\|_1-f_t(x)\}.
\]
Partitioning cells by nearest-node boundaries makes the expression separately affine in each coordinate. Its maximum is therefore at a partition vertex; the implementation encloses all such vertices. Both uniform and curvature-refined FVI receive this sharper one-sided account, not an artificially doubled baseline. Their acquired-state allowance is derived separately in \eqref{eq:fvisensor49}.

\section{Original-optimum accuracy in the investment economy}\label{sec:accuracy62}
The exact witness theorem answers an implementation question conditional on a continuation. An economist also needs to know how close the complete implemented policy is to the original optimum. We now establish that connection in the investment economy itself. The argument exploits a property not used by the earlier cell-range calculation: despite the bilinear transition, the Bellman objective remains jointly convex. This permits a second-order interpolation account and transfers verified nodal decisions to an implementable policy over the entire state domain. It does not assume that a fitted neural continuation is convex.

\paragraph{Relation to the preceding two-control query.}
The earlier query experiment used the fixed shared capacity $1/8$. The complete dynamic extension below uses the original state-dependent capacity $c(x)$, with the retained exposure columns, action-cost interaction and second shock. Its nesting claim concerns the original scalar economic model; it is not an assertion that a full dynamic policy was already present in the earlier isolated query.

\subsection{Primitives and Bellman regularity}
Let $d\ge2$ be even, let indices be cyclic, and set $\beta=15/16$. Write
\begin{align*}
 S(x)&=\frac2d\sum_j(x_j-5/8)^2+
       \frac1{4d}\sum_j(x_j-x_{j+1})^2+
       2\left(\frac12-\frac2d\sum_jx_j\right)_+^2,\\
 g(x)&=S(x)+\frac2d\sum_j(x_j-5/8)^2,\qquad
 c(x)=\frac18+\frac1{8d}\sum_jx_j.
\end{align*}
The original scalar control has $0\le a\le c(x)$, running cost $S(x)+a^2+4a^4$, and transition
\[
 F_j(x,a,z)=\frac1{16}+\frac{x_j}{2}+\frac{x_{j+1}}8+
       \frac{x_j(1-x_{j+1})}{16}+B_ja+s_jz,
\]
where $(B_j,s_j)=(1/2,1)$ at alternating coordinates and $(1/4,-1)$ at the others, and $z$ is uniform on $[-1/32,1/32]$.

For the retained two-control extension, $a_1,a_2\ge0$, $a_1+a_2\le c(x)$, the second exposure column interchanges $1/2$ and $1/4$, and the action cost is
\[
 R(a)=a_1^2+a_2^2+4(a_1^4+a_2^4)+a_1a_2/4.
\]
An independent common-direction uniform innovation of radius $1/64$ is added. This is the previously specified extension, nested at $a_2=0$ when that innovation is disabled. Both models preserve $[0,1]^d$. Let $m\in\{1,2\}$ denote the number of controls and define
\begin{align*}
 Q_t^*(x,a)&=S(x)+R(a)+\beta\mathbb E V_{t+1}^*(F(x,a,Z)),\\
 V_t^*(x)&=\min_{a\in A(x)}Q_t^*(x,a),\qquad V_T^*=g.
\end{align*}
All costs and laws in this display are the economic primitives, not a fitted surrogate.

\begin{theorem}[Regularity of the original Bellman problem]\label{thm:regularity62}
For either control specification and every finite horizon, $V_t^*$ is convex and continuously differentiable in the interior, with the corresponding one-sided extension on the cube. It is $G_t$-Lipschitz in the $\ell^1$ norm and $M_t$-semiconcave in the Euclidean norm, where valid constants are
\begin{align*}
 G_T&=10/d,&G_t&=\frac{1037}{128d}+\beta\frac{47}{64}G_{t+1},\\
 M_T&=26/d,&M_t&=\frac{22}{d}+\frac{21}{256d}
             +\beta\left(\frac9{16}M_{t+1}+\frac18G_{t+1}\right).
\end{align*}
In particular, $G_t\le27/d$ and $M_t\le54/d$. The function $Q_t^*$ is jointly strongly convex on its convex feasible graph, with state curvature at least $107/(128d)$ and action curvature at least $7/4$.

A valid semiconcavity constant along each coordinate line is
\begin{align*}
 D_T&=\frac9d+\frac{16}{d^2},\\
 D_t&=\frac5d+\frac{16}{d^2}+\frac{21}{256d^2}
 +\beta M_{t+1}\left(\frac{85}{256}+\frac{23}{256d}\right).
\end{align*}
Consequently, for the positive multilinear interpolant $I_n$ on a mesh of spacing $1/n$,
\begin{equation}\label{eq:interpolation62}
 0\le I_nV_t^*(x)-V_t^*(x)\le\kappa_t,
 \qquad \kappa_t=\frac{dD_t}{8n^2}.
\end{equation}
\end{theorem}

\noindent The proof is in Supplement, Section~\ref{supp:proof-accuracy62-print-1}.


The theorem is a structural statement about these economic primitives. It is not a claim that every controlled diffusion, recursive preference specification, or game in the complete development edition has a convex Bellman problem. Nor does it require a signed-weight neural fit to preserve convexity. Convex approximation of dynamic programs has an established literature; the additional work here is to verify the needed curvature for the original nonlinear capacity-constrained model and connect it to the actual NBO actor and complete error account.

\subsection{Certified nodal Bellman bounds}
Suppose nodal arrays satisfy $l_{t+1}\le V_{t+1}^*\le u_{t+1}$ at all grid vertices. Equation~\eqref{eq:interpolation62} gives the functions $I_nl_{t+1}-\kappa_{t+1}$ and $I_nu_{t+1}$ as valid lower and upper continuations. The lower function need not itself be convex; convexity is used for the unknown true value and for the comparison argument, not asserted of arbitrary endpoint arrays.

Divide each uniform innovation component into $q$ equal bins and use its conditional mean. Convexity gives a lower quadrature bound for the true continuation. Semiconcavity gives an upper error
\begin{equation}\label{eq:quadrature62}
 \zeta_{t+1}=\frac{M_{t+1}}{6q^2}
                      \sum_{r=1}^{s}\rho_r^2\|s_r\|_2^2,
\end{equation}
where $s=1$ or $2$, $\rho_1=1/32$, and $\rho_2=1/64$ in the extension. The expectation remains the original continuous law. At the terminal date $g$ is evaluated directly, so no terminal interpolation allowance is needed.

At a vertex $v$, let $\mathcal A_A(v)$ be the capacity simplex lattice with spacing $c(v)/A$, including all its faces. For $m=1,2$, a minimizer can be represented by neighboring lattice vertices whose coordinate variances sum to at most $m c(v)^2/(4A^2)$. A valid action smoothness constant is
\[
 H_t=\frac{21}{4}+\beta M_{t+1}\|B\|_2^2,
 \quad \|B\|_2^2=\frac{5d}{32}\ (m=1),\quad
 \|B\|_2^2=\frac{9d}{32}\ (m=2).
\]
Taylor's inequality and that barycentric representation imply
\begin{equation}\label{eq:action-cover62}
 0\le\min_{a\in\mathcal A_A(v)}Q_t^*(v,a)-V_t^*(v)
       \le\delta_t:=\frac{mH_t}{128A^2}.
\end{equation}
The linear Taylor terms cancel even when the optimum lies on a face. This is an error against the continuous feasible optimum, not only against the finite action menu.

Let $q_t^-(v,a)$ and $q_t^+(v,a)$ include all interpolation, quadrature and arithmetic errors and enclose $Q_t^*(v,a)$. Then
\begin{equation}\label{eq:nodal62}
 l_t(v)=\max\{0,\min_{a\in\mathcal A_A(v)}q_t^-(v,a)-\delta_t\},
 \qquad
 u_t(v)=\min_{a\in\mathcal A_A(v)}q_t^+(v,a)
\end{equation}
are valid nodal Bellman bounds. The zero truncation uses the nonnegative economic costs. Extra neural or conventional proposals may improve the upper side but do not remove the lower action-cover allowance. Positive interpolation, averaging and minimum maps are monotone and nonexpansive; the separate allowances in \eqref{eq:interpolation62}--\eqref{eq:action-cover62} are therefore propagated, rather than being discarded because an endpoint calculation contains the truth.

\subsection{From a fitted witness to the complete implemented policy}
Let $a_t(v)\in A(v)$ be the actual action returned at each vertex by any generator. It may be selected by the exact neural witness, a conventional fitted continuation, or their common verified fallback. Write $w_v(x)$ for the multilinear weights and set $\bar\pi_t(x)=\sum_vw_v(x)a_t(v)$. Affineness of capacity makes this action feasible. Let $q_t^{\mathrm{sel},+}(v)$ be a verified upper bound on its original-optimum continuation value $Q_t^*(v,a_t(v))$, using the same lower and upper Bellman reference for all generators.

\begin{theorem}[Actual-policy, whole-domain certificate]\label{thm:policy62}
Set
\[
 e_t=\max_v\{q_t^{\mathrm{sel},+}(v)-l_t(v)\}+\kappa_t.
\]
Suppose the implemented policy $\widehat\pi_t$ is feasible and differs from $\bar\pi_t$ by at most $b_t$ in action $\ell^1$ norm, uniformly over the state cube. Define
\[
 \Lambda_t=\frac{13}{16}+\beta\frac{3d}{8}G_{t+1},\qquad
 D_T^\pi=0,\quad D_t^\pi=e_t+\Lambda_tb_t+\beta D_{t+1}^\pi.
\]
Then, at every true state and date,
\[
 0\le J_t^{\widehat\pi}(x)-V_t^*(x)\le D_t^\pi.
\]
This statement is conditional on the actual recorded policy, not on success of its fitting algorithm. It applies equally to a pure fitted policy and to a policy augmented by verified common candidates.
\end{theorem}

\noindent The proof is in Supplement, Section~\ref{supp:proof-accuracy62-print-2}.


In the execution, observation is rounded down at forty bits. An acquired state $y\le x$ has no greater capacity than the true state. Each interpolated action is then rounded down at twenty bits using the lower endpoint of its arithmetic enclosure. If $L_{ij}$ bounds the derivative of action coordinate $i$ with respect to state coordinate $j$, and $\eta$ bounds native interpolation error, a valid implementation allowance is
\[
 b_t\le 2^{-40}\sum_{i,j}L_{ij}
           +m(2^{-20}+2\eta+\eta_{\mathrm{out}}).
\]
Here $L_{ij}$ is computed from the maximum adjacent nodal difference times $n$, and $\eta_{\mathrm{out}}$ covers endpoint subtraction. All these numbers are recorded outward. This new barycentric acquired actor is explicitly distinguished from the earlier discontinuous cell-constant actor; an old actor is not silently replaced during a purported fixed-policy revalidation.

\subsection{Arithmetic, work, and scope}
The native interpolation kernel forms nonnegative tensor weights and a weighted sum of at most $2^d$ stored binary64 endpoints. It is compiled without reassociation or floating-point contraction. For $K=(3d+4)2^d+32$, an explicit conservative absolute allowance is
\[
 \eta=\frac{K2^{-53}}{1-K2^{-53}}\max_v|z_v|+K2^{-1022}.
\]
The final term conservatively covers underflow. Input coordinates, affine indexing, finite endpoints and the binary64 format are checked. Roundoff expansion and the ordinary interval cost arithmetic remain separate from the analytic approximation bounds. Exact rational fixtures compare the native sum with its mathematical tensor interpolant.

With $N=(n+1)^d$ vertices, $K_A=A+1$ scalar actions or $K_A=(A+1)(A+2)/2$ two-control actions, the reference calculation uses $O(TNK_Aq^s2^d d)$ elementary work under the explicit tensor implementation. Stored reference values use $O(TN)$ words and a policy uses $O(TNm)$ words. Streaming batches avoid a state-pair table. Proposal fitting, exact rational witness work, actor transfer, serialization, independent cost inference and compilation are additional costs, not absorbed into this bound. The analytic errors are second order in state, action and innovation resolution; the state-count and tensor-evaluation factors still depend exponentially on dimension. The result supplies a matched original-optimum comparison without asserting that these costs have disappeared.

\section{Complete-policy qualification and endogenous occupation}\label{sec:recursive86}
The first-date result holds the continuation policy fixed. A complete policy changes both later actions and the states at which they are taken. We now retain both changes. This extension concerns the original transition kernels on the continuous state space, rather than a replacement finite-state economy.

Fix a feasible reference policy $\tau$ and its bounded original-law values $V_t^\tau$. In the expectation model of Section~\ref{sec:core49}, define
\begin{equation}\label{eq:advantage86}
 A_t^\tau(x,a)=c_t(x,a)+\beta_tP_t^aV_{t+1}^\tau(x)-V_t^\tau(x).
\end{equation}
Here $a$ can denote a whole-cell feasible action rule, including acquisition and repair. At every date let $\mathcal C_t$ be a finite measurable acquisition partition. A menu label $a\in\mathcal A_{tC}$ denotes a measurable rule $x\mapsto a(x)\in A_t(x)$ for all $x\in C$. A parameter vector $\theta\in\Theta$ specifies the menu choice in every cell at every date; ties follow the declared measurable convention. The vector may share coordinates across dates and regimes. Let $\Theta_\epsilon=\{\theta\in\Theta:F(\theta)\le U+\epsilon\}$, where $U=F(\widehat\theta)$ for a feasible attained vector. The qualification functional may use the candidate's own future. Convexity is not required for the transfer; a claimed exclusion must nevertheless have a valid certificate for this actual functional.

Write $v=(t,C,a)$ for a menu event. Retain any family $\mathcal H_\epsilon$ of finite event sets for which a verified calculation proves
\begin{equation}\label{eq:conflict86}
 \Theta_\epsilon\cap\bigcap_{v\in K}\{\theta:\pi_\theta\text{ selects }v\}=\varnothing,
 \qquad K\in\mathcal H_\epsilon.
\end{equation}
A numerical infeasibility status alone is not such a certificate. The full-menu support duals in Section~\ref{sec:compatible84} give one implementation when labels and features yield the stated linear constraints. With nonlinear own-future dependence, those particular duals cannot simply be reused; verified outer relaxations or direct interval exclusion are required. Only strict separation excludes an event, so threshold equality remains included.

For every menu event and successor cell $D$, obtain uniform original-law enclosures
\begin{align}\label{eq:occupation-input86}
 \ell_{tCa}&\le A_t^\tau(x,a(x))\le u_{tCa}, &&x\in C,\notag\\
 \underline p_{tCaD}&\le P_t(D\mid x,a(x))\le\overline p_{tCaD}, &&x\in C.
\end{align}
The bounds must cover acquisition boundaries, feasible repairs and numerical error. Initial cell probabilities $\mu_C$ are known here; certified initial-probability constraints can instead be included in the feasible set. Averages under the reference policy are not substitutes for the uniform bounds in \eqref{eq:occupation-input86}.

Define $\mathcal P_\epsilon$ by nonnegative occupation variables $m_{tC}$, $q_{tCa}$, transition flows $z_{tCaD}$, and $0\le y_{tCa}\le1$, subject to
\begin{align}\label{eq:occupation-polytope86}
 m_{0C}&=\mu_C,& \sum_a q_{tCa}&=m_{tC},& \sum_a y_{tCa}&=1,\notag\\
 q_{tCa}&\le y_{tCa},& \sum_Dz_{tCaD}&=q_{tCa},&
 m_{t+1,D}&=\sum_{C,a}z_{tCaD},\notag\\
 \underline p_{tCaD}q_{tCa}&\le z_{tCaD}\le\overline p_{tCaD}q_{tCa},&&&\notag\\[-2pt]
 \sum_{v\in K}y_v&\le |K|-1 &&(K\in\mathcal H_\epsilon).&&
\end{align}
Every $\mu_C$ is nonnegative and $\sum_C\mu_C=1$. Thus each $m_t$ is a probability vector. These are outer occupation constraints, not a claim that the acquired cell process is Markov. In particular, they allow the within-cell state law to depend on all earlier actions.

\begin{theorem}[Recursive common-parameter cost transfer]\label{thm:occupation86}
Under the preceding boundedness, feasibility and enclosure conditions, put
\begin{equation}\label{eq:recursive-band86}
 L_\epsilon=\min_{\mathcal P_\epsilon}\sum_{t,C,a}B_t\ell_{tCa}q_{tCa},\qquad
 U_\epsilon=\max_{\mathcal P_\epsilon}\sum_{t,C,a}B_tu_{tCa}q_{tCa}.
\end{equation}
Every $\theta\in\Theta_\epsilon$ satisfies
\[
 L_\epsilon\le J_0^{\pi_\theta}-J_0^\tau\le U_\epsilon.
\]
The conclusion allows changes in every date's action and continuation, together with their induced original-law occupations. Valid additional compatibility or transition constraints can only tighten the enclosure. If $\epsilon$-dependent exclusions are retained monotonically, the enclosures are nested as qualification is strengthened. Certified dual lower and upper bounds suffice; an optimizer's reported optimum does not.

For two complete-policy classes with reference-cost interval $[l_{jk},u_{jk}]$ on a simultaneous event, their cost difference is enclosed on the same event by
\[
 [\,l_{jk}+L_j-U_k,\ u_{jk}+U_j-L_k\,].
\]
No probability union over the parameter vectors is added after that event. A qualified bound on a reference policy's loss against $V_0^*$ can also be enlarged by $U_\epsilon$, retaining the original optimum as the benchmark.
\end{theorem}

The proof in Supplement~\ref{supp:occupation86} uses a telescoping identity under the candidate policy's occupation law, then embeds each complete candidate in \eqref{eq:occupation-polytope86}. Independent choices at different cells and dates are not treated as jointly attainable. As in the first-date assignment bound, the relaxation itself can have unattainable points. The economic conclusion is an enclosure, not a worst-policy realization claim.

\subsection{Nonlinear recursion, temporal selves and games}\label{sec:nonlinear86}
For a monetary evaluation $\rho_t$ satisfying the monotonicity and cash-invariance conditions in Section~\ref{sec:core49}, let $\mathcal T_t^\pi$ be the fixed-policy operator. Suppose, uniformly over the states and policies in the declared class,
\[
 d_t^-\le\mathcal T_t^\pi V_{t+1}^\tau-V_t^\tau\le d_t^+.
\]
There is no linear-occupation identity for an arbitrary nonlinear $\rho_t$. Its correct counterpart is the following recursive order statement.

\begin{proposition}[Recursive order and deviation bounds]\label{prop:nonlinear86}
If $e_T^-=e_T^+=0$ and $e_t^\pm=d_t^\pm+\beta_te_{t+1}^\pm$, then
$e_t^-\le V_t^\pi-V_t^\tau\le e_t^+$ for every class member. More generally, on invariant recursive-value domains, if a monotone fixed-policy aggregator is $\lambda_t$-Lipschitz in its continuation argument and its reference residual is at most $b_t$ in absolute value, then
\[
 \|V_0^\pi-V_0^\tau\|_\infty
 \le\sum_{t<T}\left(\prod_{j<t}\lambda_j\right)b_t.
\]
For a dynamic game, apply the one-sided sandwich of Proposition~\ref{prop:sandwich49} to each player's best-response problem with the other players' complete strategies fixed. If that account is at most $r_{it}$ in every subgame, the attained profile has unilateral cost-improvement incentive at most $r_{it}$ there. This is an approximate equilibrium certificate, not an equilibrium-existence assertion. Date-specific temporal selves are treated as distinct decision makers, with their stated objectives and future strategies, rather than as one artificially time-consistent planner.
\end{proposition}

Endogenous preference states belong in the augmented state and transition before these hypotheses are checked. A diffusion approximation contributes its certified one-step operator error at every date, propagated with the same continuation moduli; an unproved weak-discretization rate is not assumed. The broader economic applications and their original hypotheses remain intact. The expectation identity, nonlinear order bound and strategic deviation certificate are distinct results, not interchangeable interpretations of one simulation table.

\section{A common operator beyond the observed catalogue}\label{sec:population86}
A coefficient vector checked at finitely many tasks need not work at a new task. The constructive NBO has a stronger route when the economic primitives provide joint regularity. We state the necessary regularity on states and contexts rather than assuming successful extrapolation by a fitted nine-coefficient readout.

Let $r=(s,x)$ range over an invariant compact subset $K_t$ of $\R^{p_t}$, where $s$ contains declared economic contexts and any endogenous preference variables. Actions and transitions are defined on this augmented domain. Suppose the primitive moduli, Hausdorff-continuous feasible menus and measurable repairs in Theorem~\ref{thm:construct49} hold jointly in $r$, with constants $L_t,D_t,K_t^A$. Let an effective cover have $S_t$ nodes and radius $h_t$. At each node obtain its own-future minimized label with lower error $-e_t$ and upper error $D_tk_t+e_t$. Repair displacement is bounded by $K_t^A\|r-r_i\|_1+\zeta_t$.

\begin{theorem}[Population construction with one shared neural circuit]\label{thm:population86}
The single function
\[
 f_t(r)=\min_{i\le S_t}\{y_{ti}+L_t\|r-r_{ti}\|_1\}
\]
has an exact affine--ReLU circuit, a measurable original-owner actor and the bounds \eqref{eq:enclosures49} uniformly on $K_t$. Consequently \eqref{eq:constructionloss49}, with the joint-domain constants, holds for every initial context and every admissible evolution remaining in those domains. Contexts may be generated endogenously and adaptively. No independence assumption on their deployment sequence is needed for this uniform deterministic assertion.

The circuit uses $O(p_tS_t)$ nonzero scalar operations and weights, with absolute-value gates and a balanced minimum tree of depth $O(\log S_t)$. If the joint domains have effective $\ell^1$ covering numbers at most $C_th^{-k_t}$, this yields $S_t\le C_th_t^{-k_t}$. It is the explicitly shared joint-domain circuit, not a collection of task-specific networks. The original-owner and repair obligations remain necessary.
\end{theorem}

This is a constructive population theorem, not a claim that the small feature maps in the historical experiments satisfy its approximation bound. The covering dimension $k_t$, the ambient readout cost $p_t$, acquisition, integration and label work are charged separately. A low intrinsic covering dimension must be established for the declared economic family. The same cone also has a native min-plus implementation: exact equivalence of two encodings is not neural superiority.

\begin{proposition}[Acquisition coverage under persistent exploration]\label{prop:coverage86}
Choose $M$ deterministic radius-$h$ balls whose centers cover $K$. If, before each acquisition $i$, every such ball has conditional probability at least $a_h>0$ of containing the next node, then after $n$ acquisitions the probability that the acquired nodes fail to form a radius-$2h$ cover is at most $M(1-a_h)^n\le M\exp(-na_h)$. Thus $n\ge[\log M+\log(1/\delta)]/a_h$ is sufficient. The assertion permits an adaptive acquisition law satisfying the conditional lower bound; it does not silently replace such a law by IID sampling. On the coverage event, Theorem~\ref{thm:population86} applies with $h_t=2h$, provided the own-future labels and primitive moduli have their stated simultaneous certification.
\end{proposition}

\subsection{From approximation to predictor work}
The certificate's validity and a proposal's value are different. Consider a complete finite action menu at a query, with true objective $Q$ and estimates having simultaneous radius at most $e$. Select a minimum estimated objective and close only when its upper endpoint does not exceed every competitor's lower endpoint. Let $\Delta$ be the true best-to-second-best gap, with $\Delta=0$ for ties.

\begin{proposition}[A distributional work consequence of uniform accuracy]\label{prop:work86}
The preceding test closes whenever $\Delta>4e$. Suppose, uniformly over the relevant policy-induced query laws, the tie probability is at most $p_0$ and $\Pr(0<\Delta\le z)\le Cz^\alpha$ for $z>0$. If proposal and certification work are at most $h(e)$ and a protected fallback uses at most $W_+$, then
\[
 \E W\le h(e)+W_+\min\{1,p_0+C(4e)^\alpha\}.
\]
If the always-conventional comparator has expected work at least $W_-$, a positive saving follows when the displayed upper bound is below $W_-$. Comparing two upper bounds is not sufficient. The assertion is componentwise for nonnegative operation ledgers, or for a separately declared nonnegative resource-price vector. It is not a wall-clock assertion without a valid mapping from the ledger to time.
\end{proposition}

Uniformity over induced query laws is essential: a margin observed under a baseline state distribution does not establish the bound after the policy changes. The tie term cannot be discarded. A representation-specific empirical comparison must measure $h(e)$, achieved accuracy and the relevant query distribution for each representation. The constructed neural circuit provides an explicit member of this comparison; protected fallback alone supplies no evidence about neural training reliability. Proofs are in Supplement~\ref{supp:population86}.

\section{Signed transition qualification for recursive preferences}\label{sec:risk91}
The transport account can use more information than an absolute error radius. A numerical transition may have a known direction of bias, and the implemented action may have a more precise upper error bound than the whole feasible menu. We retain these distinctions. The resulting qualification applies to the original recursive evaluation at every date, not only to expected cost with a fixed continuation. It does not change an expectation objective into a risk objective or identify a static tail-risk criterion with its nested counterpart.

Retain the augmented economic state, original feasible actions and measurable complete policies of Section~\ref{sec:core49}. Suppose bounded original values satisfy backward recursion with
\[
 Q_t(v;x,a)=c_t(x,a)+\beta_t\rho_{t,x,a}(v),\qquad
 T_tv(x)=\inf_{a\in A_t(x)}Q_t(v;x,a),
\]
where each original conditional evaluation is monotone and cash invariant: $\rho(v+k)=\rho(v)+k$ for constants $k$. Discounts $\beta_t\ge0$ are deterministic, $B_0=1$, and $B_t=\prod_{s<t}\beta_s$. The state includes the preference variables required by the original model. The relevant bounded functions are in the evaluation's domain, and the original Bellman optimum and supplied feasible policies are well defined. For each common parameter $\theta$, store bounded $f_{t,\theta}$, a complete feasible $\pi_\theta$, and numerical action values $\widehat Q_{t,\theta}$ on the original feasible actions.

Put $\widehat T_{t,\theta}(x)=\inf_a\widehat Q_{t,\theta}(x,a)$ and define the signed discrepancy
\[
 D_{t,\theta}(x,a)=Q_t(f_{t+1,\theta};x,a)-\widehat Q_{t,\theta}(x,a).
\]
Assume uniform, verified bounds
\begin{align}\label{eq:signed-input91}
 \ell_{t,\theta}&\le\widehat T_{t,\theta}(x)-f_{t,\theta}(x)
                         \le u_{t,\theta},\notag\\
 0&\le\widehat Q_{t,\theta}(x,\pi_{t,\theta}(x))
                          -\widehat T_{t,\theta}(x)\le\gamma_{t,\theta},\notag\\
 b^-_{t,\theta}&\le D_{t,\theta}(x,a)\quad(a\in A_t(x)),\notag\\
 D_{t,\theta}(x,\pi_{t,\theta}(x))&\le b^{\pi,+}_{t,\theta}.
\end{align}
Here $b^-_{t,\theta}$ and $b^{\pi,+}_{t,\theta}$ may be signed. A bound confined to sampled actions cannot replace the all-action lower bound. Original-owner repairs, unsampled feasible actions, integration error and arithmetic error must be included in the relevant inequalities. Let $d^-_\theta\le g-f_{T,\theta}\le d^+_\theta$ throughout the terminal domain.

\begin{theorem}[Signed, recursive original-optimum qualification]\label{thm:signed91}
Under the stated conditions, every supplied complete policy satisfies, at every initial state,
\begin{align}\label{eq:signed-loss91}
 0\le V_0^{\pi_\theta}-V_0^*
 &\le \mathcal S_\theta,\notag\\
 \mathcal S_\theta
 &=\sum_{t<T}B_t\bigl\{u_{t,\theta}-\ell_{t,\theta}
       +\gamma_{t,\theta}+b^{\pi,+}_{t,\theta}-b^-_{t,\theta}\bigr\}\notag\\
 &\quad+B_T(d^+_\theta-d^-_\theta).
\end{align}
Consequently a class for which these inequalities hold uniformly and $\sup_\theta\mathcal S_\theta\le\varepsilon$ consists entirely of complete policies with original-recursion loss at most $\varepsilon$. Actions, attained continuations and endogenous occupations may change at every date. The assertion is pointwise in the augmented state and does not require independent deployment contexts or a continuous actor.

The bound is invariant to arbitrary date-specific, state-independent shifts of the stored continuations when the numerical action values are shifted consistently. If the original-minus-numerical action value is a constant $k_{t,\theta}$ independent of state and action, its contribution to \eqref{eq:signed-loss91} is zero. A bias common to all actions at a given state, but varying across states, does not in general cancel. The coefficient one on $b^{\pi,+}-b^-$ cannot be reduced uniformly, even for one decision date and deterministic transitions.
\end{theorem}

The proof compares both the complete policy and the original optimum with the same stored continuation by recursive order. It does not propagate a linear occupation identity through nonlinear preferences. Setting $b^-=-b$ and $b^{\pi,+}=b$ recovers Theorem~\ref{thm:transport90}; a verified signed enclosure can strictly improve it without a new fit or economic observation. With stationary bounded data and $0\le\beta<1$, the corresponding stationary bound is
\[
 V^\pi-V^*\le
 \frac{u-\ell+\gamma+b^{\pi,+}-b^-}{1-\beta}.
\]
The proof and the necessary bias counterexamples are in Supplement~\ref{supp:signed91}.

\subsection{Explicit qualification of two nonlinear evaluations}
Wasserstein distance controls expectations of Lipschitz functions. For nonlinear evaluations its price must be derived, not borrowed from the expectation model. For a bounded $K$-Lipschitz continuation $f$ and probability laws $P,\widehat P$, define the upper-tail average value at risk and entropic evaluation by
\begin{align}\label{eq:risk-definitions91}
 \rho^\alpha_P(f)&=\inf_{z\in\R}
       \left\{z+\frac{P(f-z)_+}{1-\alpha}\right\},
       &&0\le\alpha<1,\notag\\
 \rho^\lambda_P(f)&=\lambda^{-1}\log P\exp(\lambda f),
       &&\lambda>0.
\end{align}
The limit at $\lambda=0$ is expectation. The first definition includes atoms; it is not the conditional mean obtained by discarding all mass at a quantile. The variational representation and its treatment of general loss distributions are established results \citep{rockafellaruryasev2002}. We use these evaluations only when they are the original one-step primitives of the recursive economic objective.

\begin{proposition}[Risk-specific transport allowances]\label{prop:risk-transport91}
If $W_1(P,\widehat P)\le e$, then
\begin{align}\label{eq:risk-moduli91}
 |\rho^\alpha_P(f)-\rho^\alpha_{\widehat P}(f)|
       &\le \frac{K e}{1-\alpha},\notag\\
 |\rho^\lambda_P(f)-\rho^\lambda_{\widehat P}(f)|
       &\le K\exp\{\lambda\operatorname{osc}(f)\}\,e.
\end{align}
The entropic expression is invariant to constants added to $f$. The factor $(1-\alpha)^{-1}$ is necessary uniformly over bounded Lipschitz continuations and probability laws. If a coupling satisfies $f(Y)\ge f(\widehat Y)$ almost surely, the original-minus-numerical risk difference is nonnegative for either evaluation.
\end{proposition}

If the additional numerical action-value error is at most $\xi_t$ in absolute value, a symmetric choice in \eqref{eq:signed-input91} is
\[
 b^-_t=-\beta_t\kappa_t e_t-\xi_t,
 \qquad b^{\pi,+}_t=\beta_t\kappa_t e_t+\xi_t,
\]
where $\kappa_t=K_{t+1}/(1-\bar\alpha_t)$ for a verified risk-level cap $\alpha\le\bar\alpha_t<1$, or $\kappa_t=K_{t+1}\exp(\bar\lambda_t O_{t+1})$ for verified bounds on the entropic parameter and continuation oscillation. A monotone coupling instead permits $b^-_t=-\xi_t$ and the same upper endpoint. Approximate evaluation of the exponential requires an outward upper bound; machine output is not an exact modulus.

These local risk coefficients enter the dated budget additively:
\begin{equation}\label{eq:risk-budget91}
 \sum_{t<T}B_t\eta_t
 +2\sum_{t<T}B_{t+1}\kappa_t e_t
 +2\sum_{t<T}B_t\xi_t+B_T(d^+-d^-)\le\varepsilon,
\end{equation}
where $\eta_t$ bounds residual span plus greediness. Monotonicity and cash invariance propagate constant error by $\beta_t$, not by a product of all later tail-risk multipliers. The continuation moduli themselves may grow with horizon or preference parameters and must still be established and paid for. Subject to this qualification, the coupled resource program of Section~\ref{sec:coupled89} allocates residual and risk-adjusted transition precision without deleting descendant or representation costs.

\paragraph{Economic and strategic interpretation.}
A larger tail-risk level can require finer transition information for the same original-policy loss allowance. That requirement is part of the original preference primitive, not a favorable reweighting of measured work. When preference levels evolve endogenously, the bounds must cover the augmented state and actual feasible actions. For games and temporal selves, the theorem applies to each fixed-opponent or fixed-future best-response recursion under that decision maker's original objective. Applying nested average value at risk to an originally static terminal-risk problem would change that problem; no such substitution is made. General non-cash-invariant recursive aggregators retain Proposition~\ref{prop:nonlinear86} and its separate continuation modulus.

\section{Reference-centered risk kernels and complete-policy qualification}\label{sec:centeredrisk92}
A datewise worst-case error account can charge incompatible states at successive dates. Under nonlinear recursive preferences, replacing that account by an expectation under the physical transition law is not valid. We instead retain the conditional risk evaluation and center its dual representation at the actual stored NBO continuation. The resulting state-dependent certificate has an original-optimum interpretation and a finite occupation representation. It extends, rather than replaces, the signed qualification of Section~\ref{sec:risk91}.

\subsection{The centered operator}
Retain the original augmented economic state, feasible actions, nonnegative discounts, and bounded complete-policy recursions. At each $(t,x,a)$ suppose the original conditional evaluation admits the representation
\begin{equation}\label{eq:dualrisk92}
 \rho_{t,x,a}(v)=\sup_{q\in\mathcal Q_{t,x,a}}
 \left\{\int v(y)q(dy)-\psi_{t,x,a}(q)\right\}.
\end{equation}
Every $q$ is a probability measure on the original successor state space. The representation is assumed on a common bounded-function domain containing all continuations below, with the measurability and selection conditions needed for the original Bellman recursions. This is a hypothesis about the original preference primitive, not permission to replace a static risk objective by a nested one. Conditional risk duality and risk-averse dynamic programming are established principles \citep{ruszczynski2010risk,ruszczynski2014erratum}. For a nonconvex or non-cash-invariant aggregator lacking \eqref{eq:dualrisk92}, the earlier order and modulus results remain applicable.

Fix one qualified parameter $\theta$ and suppress its subscript. Its entire sequence of stored continuations is $f=(f_t)_{t=0}^T$, and its supplied complete feasible policy is $\pi$. Define, using the same $f_{t+1}$ in every term,
\begin{align}\label{eq:centered-def92}
 r_t(x,a)&=c_t(x,a)+\beta_t\rho_{t,x,a}(f_{t+1})-f_t(x),\notag\\
 k_t^f(x,a,q)&=\psi_{t,x,a}(q)+\rho_{t,x,a}(f_{t+1})
                         -\int f_{t+1}\,dq,\notag\\
 \mathcal E_{t,x,a}^f(z)&=\sup_{q\in\mathcal Q_{t,x,a}}
                   \left\{\int z\,dq-k_t^f(x,a,q)\right\}.
\end{align}
The penalty is nonnegative and has infimum zero, whether or not the supremum in \eqref{eq:dualrisk92} is attained. Direct substitution gives
\begin{equation}\label{eq:exact-center92}
 \mathcal E_{t,x,a}^f(z)
 =\rho_{t,x,a}(f_{t+1}+z)-\rho_{t,x,a}(f_{t+1}).
\end{equation}
In particular, $\mathcal E^f(0)=0$, and $\mathcal E^f$ is monotone, cash invariant, and nonexpansive in the uniform norm. Discarding $k^f$ is an upper relaxation, not this identity. Freezing one risk-maximizing kernel at $f$ is in general not even an upper relaxation.

\begin{theorem}[State-dependent original-optimum certificate]\label{thm:centeredrisk92}
Suppose verified functions satisfy $r_t^-(x,a)\le r_t(x,a)$ for every original feasible action and $r_t(x,\pi_t(x))\le r_t^{\pi,+}(x)$. Let bounded measurable $L_t,U_t$ satisfy
\begin{align}\label{eq:potential-cert92}
 L_T&\le g-f_T\le U_T,\notag\\
 L_t(x)&\le\inf_{a\in A_t(x)}
       \{r_t^-(x,a)+\beta_t\mathcal E_{t,x,a}^f(L_{t+1})\},\notag\\
 U_t(x)&\ge r_t^{\pi,+}(x)
                +\beta_t\mathcal E_{t,x,\pi_t(x)}^f(U_{t+1}).
\end{align}
Then, at every initial state,
\begin{equation}\label{eq:potential-loss92}
 f_0+L_0\le V_0^*\le V_0^\pi\le f_0+U_0,
 \qquad 0\le V_0^\pi-V_0^*\le U_0-L_0.
\end{equation}
The same statements hold at every intervening date and state. They are uniform over a qualified parameter class whenever all premises and the reported bound hold for each member with its own common continuation sequence and complete policy. The certificate permits endogenous states and actions changing at every date; it does not require a finite state space or independent deployment contexts.

With exact residuals and terminal discrepancies, equality recursions in \eqref{eq:potential-cert92} give $L_t=V_t^*-f_t$ and $U_t=V_t^\pi-f_t$. With scalar bounds from \eqref{eq:signed-input91}, their corresponding constant lower and upper sequences are feasible certificates here. Thus the exact state-dependent recursion is no looser than the signed scalar account. A strictly smaller bound can result without changing a fit, a feasible action, or an economic observation.
\end{theorem}

The last assertion concerns the best exact recursion or a verified improvement upon the constant certificates. It is not a claim that every approximately computed $L,U$ improves the scalar account. In implementation, the minimum of valid policy supersolutions and the maximum of valid optimum subsolutions retain validity, so an inferior new bound need not displace the existing one.

The lower inequality requires every original action. A discrete action search without a verified original-menu lower bound cannot establish it. For a fixed action, a feasible risk kernel supplies a lower bound on $\mathcal E^f(L)$; an upper bound requires control of the entire risk envelope. Numerical integration, feasible repairs, residual errors and penalty errors belong in these inequalities. If penalty enclosures are used, lower penalty endpoints enter an upper certificate and upper endpoints enter a lower witness. Approximate dual optimizers are accepted only through these independently checked inequalities.

\subsection{Exact risk occupations and safe kernel exclusion}
Consider now a finite-state implementation for which the original one-step risk has a finite affine representation
\[
 \rho_{t,s,a}(v)=\max_{j\in\mathcal J_{t,s,a}}
       \{q_{t,s,a,j}^{\mathsf T}v-\psi_{t,s,a,j}\}.
\]
The supplied complete policy is fixed before the following maximization. Write $r_{ts}=r_t(s,\pi_t(s))$, $k_{tsj}=k_t^f(s,\pi_t(s),q_j)$ and $q_{tsj}=q_{t,s,\pi_t(s),j}$. For an initial probability vector $\mu$, let $\mathcal Y^\pi$ consist of nonnegative masses satisfying
\begin{equation}\label{eq:riskflow92}
 m_{0s}=\mu_s,\qquad \sum_j y_{tsj}=m_{ts},\qquad
 m_{t+1,z}=\sum_{s,j}y_{tsj}q_{tsj,z}.
\end{equation}
These are risk-distorted occupations, not frequencies under the physical law. Initial integration by $\mu$ is outside the conditional value function; a different original initial risk aggregation needs its own evaluation.

\begin{proposition}[Penalized occupation identity and dual witness]\label{prop:riskflow92}
For the preceding finite representation,
\begin{align}\label{eq:riskflow-value92}
 \mu^{\mathsf T}(V_0^\pi-f_0)
 =\max_{\mathcal Y^\pi}\biggl\{
 &\sum_{t<T,s}B_t r_{ts}m_{ts}
 -\sum_{t<T,s,j}B_{t+1}k_{tsj}y_{tsj}\notag\\
 &+B_T\sum_s(g_s-f_{Ts})m_{Ts}\biggr\}.
\end{align}
Any feasible occupation gives a lower bound on this fixed-policy value. Any $U$ with $U_T\ge g-f_T$ and
\[
 U_{ts}\ge r_{ts}+\beta_t(q_{tsj}^{\mathsf T}U_{t+1}-k_{tsj})
 \quad\hbox{for every }t,s,j
\]
gives an upper bound $\mu^{\mathsf T}U_0$. Matching witnesses certify equality using rational arithmetic when the input is rational. The optimal-policy recursion retains the order $\inf_a\sup_q$; it is not obtained by minimizing jointly over policy actions and risk kernels.
\end{proposition}

A stored continuation can also exclude risk kernels, not just action events. Fix a feasible retained witness $q_0$ and a verified set $\mathcal D$ containing the successor error vector to be evaluated. Define its support function $h_{\mathcal D}(d)=\sup_{z\in\mathcal D}d^{\mathsf T}z$. Then the strict certificate
\begin{equation}\label{eq:kernel-cut92}
 h_{\mathcal D}(q-q_0)<k^f(q)-k^f(q_0)
\end{equation}
excludes $q$ from the maximizing affine pieces throughout $\mathcal D$. Equality is retained. A verified upper support bound and lower bound on the penalty difference suffice. For a box $L\le z\le U$ the support is
\[
 \sum_i(q_i-q_{0i})_+U_i
       +\sum_i\min\{q_i-q_{0i},0\}L_i.
\]
For $\operatorname{osc}(z)\le s$, it is $s\|q-q_0\|_1/2$. The error set must be justified before exclusions are used; defining it from the same unverified pruned computation is circular. Unpruned signed or state-dependent certificates supply one valid starting set.

\subsection{Preference primitives and computational implications}
For average value at risk, $q_i\le p_i/(1-\alpha)$, $q_i\ge0$ and $\sum_iq_i=1$, with zero penalty on that capped simplex. Hence $k^f(q)=\rho_p^\alpha(f)-q^{\mathsf T}f$. Quantile atoms and zero-probability outcomes are retained \citep{rockafellaruryasev2002}. For entropic evaluation, put
\[
 p_f(dy)=\frac{e^{\lambda f(y)}p(dy)}{\int e^{\lambda f}\,dp}.
\]
Its centered penalty is $k^f(q)=\lambda^{-1}{\rm KL}(q\Vert p_f)$, and therefore
\begin{equation}\label{eq:entropic-center92}
 \mathcal E^f(z)=\lambda^{-1}\log\int e^{\lambda z}\,dp_f.
\end{equation}
The change of measure is a computational identity, not a change to the original economic law or its risk parameter. Supplement~\ref{supp:centeredrisk92} proves the representation and its support conditions. A finite affine occupation program is exact only when the stated finite representation is exact; discretizing the entropic dual without an enclosure does not meet that condition.

For a fixed finite representation, backward certificate evaluation uses at most a constant multiple of
\[
 \sum_{t,s,a}\sum_{j\in\mathcal J_{t,s,a}}
       |\operatorname{supp}(q_{t,s,a,j})|
\]
arithmetic operations once residuals and penalties are available. It stores successive value layers rather than expanding descendant histories. Kernel construction, original-law integration, neural-circuit evaluation, support certificates, rational bit lengths and stored witnesses are additional costs. Risk-envelope enumeration can itself be large. This bound does not remove state-dimension dependence, certify an unverified continuous-state aggregation, or prove a speed advantage over another implementation.

For fixed penalties, increasing residuals by nonnegative local allowances $e_{ts}$ increases the fixed-policy upper value by no more than
\begin{equation}\label{eq:occupation-budget92}
 \max_{\mathcal Y^\pi}\sum_{t,s}B_t e_{ts}m_{ts},
\end{equation}
including a terminal allowance at $t=T$. This gives an occupation-sensitive sufficient precision budget in place of $\sum_t B_t\max_s e_{ts}$. The budget prices allowable numerical errors; an application-specific price of computation remains a separate input. Shared-parameter restrictions must be retained before taking an envelope over policies, as in Section~\ref{sec:pattern89}.

\paragraph{Interpretation.}
The new mechanism depends on the stored continuation through $k^f$ and on the original preference through its risk envelope. Equivalent neural and non-neural representations have identical centered guarantees. A representation-specific computational benefit requires a matched execution, not a validity theorem. The existing held-out ReLU result, warning records, unresolved economic comparisons, and adverse vector horizon evidence therefore remain unchanged. The exact finite diagnostics below test the nonlinear certificate and its independent implementation; they are not a substitute for the original large-vector comparison or new economic observations.

\section{Continuous-law qualification in the investment model}\label{sec:continuous95}
The risk theorem becomes a numerical method only when its inequalities can be constructed from economic primitives. We now supply this construction for the original investment technology with a nested tail-risk preference. States and innovations remain continuous, every feasible action is covered, and the implemented terminal controller rather than a hypothetical exact optimizer determines policy cost. This application joins continuation construction, the original-optimum comparison and a representation comparison on one problem. It does not replace the general NBO construction or the other economic applications.

\subsection{Primitives and the object being approximated}
Let $d\geq2$ be even, $X=[0,1]^d$, and interpret subscripts cyclically. At dates $t=0,1$, the feasible investment is $0\leq a\leq A(x)=1/8+\bar x/8$, where $\bar x=d^{-1}\sum_jx_j$. Write $b_j=1/2$ for odd $j$ and $b_j=1/4$ for even $j$, and $s_j=(-1)^{j+1}$. The original transition is
\begin{equation}\label{eq:investment95}
 F_j(x,a,z)=\frac1{16}+\frac{x_j}{2}+\frac{x_{j+1}}8
       +\frac{x_j(1-x_{j+1})}{16}+b_ja+s_jz,
 \qquad z\sim\operatorname{Unif}[-1/32,1/32].
\end{equation}
Define
\begin{align}\label{eq:cost95}
 S(x)&=\frac2d\sum_j(x_j-5/8)^2+
       \frac1{4d}\sum_j(x_j-x_{j+1})^2+2(1/2-2\bar x)_+^2,\\
 g(x)&=S(x)+\frac2d\sum_j(x_j-5/8)^2,\qquad h(a)=a^2+4a^4.\nonumber
\end{align}
We use $\beta=15/16$ and, separately, $\alpha\in\{1/2,3/4\}$. Average value at risk is the upper-tail cost functional
\[
 \rho_\alpha(v)=\sup\left\{\int vq\,dP:
             0\leq q\leq(1-\alpha)^{-1},\ \int q\,dP=1\right\}.
\]
Here $P$ is the original uniform innovation law. The objective is recursively defined by $V_2=g$ and
\[
 V_t(x)=S(x)+\inf_{0\leq a\leq A(x)}
       \{h(a)+\beta\rho_\alpha(V_{t+1}(F(x,a,\cdot)))\}.
\]
It is not the static tail risk of total realized cost. A complete policy value replaces each infimum by the action prescribed by that policy, with the same recursive risk primitive.

\paragraph{Assumptions for the continuous construction.}
For this subsection, the primitives are exactly \eqref{eq:investment95}--\eqref{eq:cost95}; $\alpha$ is fixed before construction; the fitted continuation $f$ is a finite convex function on $X$; and every numerical lower or upper value is an outward enclosure. Feasible tail densities and feasible actions are retained explicitly. Terminal selection is the measurable, dyadically acquired controller described below, with its fallback included in the error bound. These concrete conditions imply compact feasible graphs, bounded costs, measurable minimizers, and the conditional-risk dual used in Theorem~\ref{thm:centeredrisk92}. They do not assert that every recursive preference or every trained network has this convex structure.

\begin{proposition}[Primitive curvature and state supports]\label{prop:curvature95}
For either risk level and every even $d$, the one-date objective
$Q_1(x,a)=S(x)+h(a)+\beta\rho_\alpha(g(F(x,a,\cdot)))$
is jointly convex on $X\times[0,1/4]$, with curvature bounded below by
\begin{equation}\label{eq:curvature95}
 \frac{4-5\beta/4}{d}\|v\|_2^2+2w^2
\end{equation}
in a direction $(v,w)$. Its value $V_1$ is convex on $X$ and satisfies
\begin{equation}\label{eq:lipschitz95}
 |V_1(x)-V_1(y)|\leq\frac{15}{d}\|x-y\|_1.
\end{equation}
Let $W$ be the objective evaluated under any fixed feasible risk density, at an anchor $(x_0,a_0)$, and let $(p_x,p_a)$ be its gradient there. For any $\lambda,\nu\geq0$, set $e=p_a+\lambda-\nu$. Then the following is a global affine minorant of $V_1$:
\begin{align}\label{eq:affine95}
 \ell(y)={}&W+\lambda(a_0-A(x_0))-\nu a_0-\frac{e^2}{4}\nonumber\\
 &+\left(p_x-\frac{\lambda}{8d}\boldsymbol{1}\right)^{\!\top}(y-x_0).
\end{align}
Neither optimal multipliers nor an exact optimizing action are required.
\end{proposition}

The curvature is obtained from the original nonlinear technology, not imposed on an interpolated value grid. The negative second-order term from $x_j(1-x_{j+1})$ is bounded using $|\partial_jg|\leq10/d$; the stage curvature dominates it. The allowance $e^2/4$ pays for imperfect stationarity. The proof in the Technical Supplement also establishes the global state modulus in \eqref{eq:lipschitz95}, including the state-dependent action constraint.

\subsection{An all-action certificate under the original innovation law}
Partition the innovation interval into $n$ equal bins $B_i$, with midpoints $z_i$ and probabilities $p_i=1/n$. An innovation bin is a set over which an inequality holds, not a replacement point mass in the economic model. Choose nonnegative $q_i\leq p_i/(1-\alpha)$ summing to one. These masses define a feasible density constant on each bin. For a fixed initial state and feasible proposal $a_0$, let $W_0$ be a lower bound on the objective under this density, obtained from globally valid affine supports of the continuation at the successor midpoints. Let $p$ enclose the derivative of the same supporting objective at $a_0$. Then
\begin{equation}\label{eq:actionlower95}
 \underline V_0(x)=\inf_{0\leq u\leq A(x)}
         \{W_0+p(u-a_0)+(u-a_0)^2\}
\end{equation}
is a lower bound on the original Bellman optimum when the supported continuation is $V_1$. The scalar quadratic is minimized over the full original interval. Interval uncertainty in $p$ is covered by taking the smaller minimum at its two endpoints, since the minimum is concave in $p$.

For a convex fitted continuation, midpoint supports give the lower risk witness and the larger endpoint value bounds the continuation throughout each bin. For $V_1$, Proposition~\ref{prop:curvature95} supplies midpoint supports and the allowance $15/(32n)$ covers each half-bin. Upper tail risk is the average of the largest $n(1-\alpha)$ upper bin bounds. This average is the exact risk of a dominating step function, hence an upper bound under the continuous law.

\begin{theorem}[Constructed original-optimum qualification]\label{thm:continuous95}
Suppose the supports and bin bounds above are outwardly verified. Suppose further that the complete terminal controller satisfies
$V_1\leq V_1^{\widehat\pi}\leq V_1+\varepsilon_1$ everywhere on $X$.
For every queried initial state, the lower value \eqref{eq:actionlower95} and the binwise upper bound $\overline J_0^{\widehat\pi}$ satisfy
\begin{equation}\label{eq:originalgap95}
 0\leq J_0^{\widehat\pi}(x)-V_0(x)
       \leq\overline J_0^{\widehat\pi}(x)-\underline V_0(x).
\end{equation}
For any stored convex continuation $f$, the same lower and upper constructions are admissible bounds for the centered-risk comparison of Theorem~\ref{thm:centeredrisk92}. Their validity does not depend on the representation of $f$, on a successful fit, or on the proposal being an exact minimizer.
\end{theorem}

To see the connection, write a verified continuation bracket as $\ell\leq V_1\leq u$ and use errors $L=\ell-f$, $U=u-f$. The centered identity gives
\[
 \rho_\alpha(f)+\mathcal E^f(L)=\rho_\alpha(\ell),\qquad
 \rho_\alpha(f)+\mathcal E^f(U)=\rho_\alpha(u).
\]
Joint evaluation avoids subtracting and then adding two separately rounded copies of $\rho_\alpha(f)$. The bound \eqref{eq:actionlower95} uses one feasible risk density to minorize the required all-action infimum; the policy upper bound controls the full risk envelope. Thus the application constructs sufficient inequalities for the original theorem. It does not claim that a non-neural continuation would receive a different validity theorem, or that computing the original-value witness is free.

The terminal controller observes each state coordinate rounded down to $40$ binary places. It accepts its floating proposal only after an outward all-action gap of at most $2^{-21}$ at that observed state; otherwise it uses rational bisection. The original-state transfer is less than $30\,2^{-40}$. Consequently $\varepsilon_1=2^{-20}$ covers the complete controller, including observation and fallback. Policy evaluation encloses this controller at every possible successor, not only at the midpoint states used by the integrator.

\subsection{Residual covers and strict continuum screening}
A cellwise residual bracket need not solve another action problem at every state. On each rectangular state cell, convex interpolation of the vertex upper values dominates $V_1$; each affine function \eqref{eq:affine95} minorizes it globally. The same interpolation and subgradient construction applies to convex $f$. Comparing an affine minorant of one function with the convex interpolation upper bound of the other at every vertex proves a bound for $V_1-f$ throughout the cell. Taking the best available anchor and then the worst cell gives a whole-domain residual interval. The $32\times32$ cover below therefore qualifies the continuous square, rather than converting its $1,089$ nodes into empirical generalization observations.

Let $[l_i,u_i]$ enclose an original successor value over every point of bin $B_i$. If $\tau$ is the $k$th largest $l_i$, $k=n(1-\alpha)$, then
\begin{equation}\label{eq:strict95}
 u_i<\tau
\end{equation}
excludes all risk mass from bin $i$ in an upper-tail maximizer. At least $1-\alpha$ of the original probability mass has value at least $\tau$; strict separation makes any positive mass on the excluded bin suboptimal. Equality is retained. These bounds must be available before screening. The screening calculation below pays for their construction and for the test itself, and compares identical upper-tail values with and without screening. This continuum-bin result is separate from enumerating all extreme kernels of a finite risk envelope.

\subsection{A fixed-accuracy horizon account for the original scalar and vector economies}\label{sec:horizon95}
The continuous two-date implementation is not the limit of the constructive argument. The original scalar model and its retained two-control, two-innovation extension in Section~\ref{sec:accuracy62} admit a horizon-uniform primitive account under nested AVaR. For the vector extension, actions satisfy $a_1,a_2\geq0$, $a_1+a_2\leq A(x)$; the second exposure column interchanges $1/2$ and $1/4$, the action cost is $a_1^2+a_2^2+4(a_1^4+a_2^4)+a_1a_2/4$, and an independent common-direction shock of radius $1/64$ is added. These are the retained original vector primitives, not a new model selected to obtain a horizon bound.

\begin{proposition}[Uniform risk regularity and constructive accuracy]\label{prop:horizon95}
Consider either economy at any finite horizon, with $\beta=15/16$ and a prespecified risk level $\alpha_t\in\{1/2,3/4\}$ at each date. Then the original value functions are convex and $27/d$-Lipschitz in the $\ell^1$ norm. Their Bellman objectives have joint state curvature at least $107/(128d)$ and action curvature at least $7/4$ on the feasible graph. No continuous differentiability of a risk-averse value function is asserted.

The own-future cone construction of Theorem~\ref{thm:construct49} may use the same modulus $27/d$ at every date. Let its state grid have $n_x$ intervals per coordinate; let the action net have resolution $m$ on the original capacity interval or capacity simplex; and let each innovation coordinate have $n_z$ equal bins, where $n_z$ is a multiple of four. With exact labels for the finite-bin risk calculation, the charged continuous-law integration allowance and exact feasible repairs, the constructed complete policy satisfies
\begin{equation}\label{eq:uniformhorizon95}
 J_0^\pi-V_0^*\leq16\left\{\frac{27}{n_x}
             +\frac{1319}{512m}+\frac{1215}{512n_z}\right\}.
\end{equation}
The bound holds for every horizon with these same resolutions. Acquisition and final action-rounding allowances are added according to \eqref{eq:acquire49}; reserving any positive share of the target tolerance for them gives finite sensor and action precisions independent of the horizon.
\end{proposition}

For an approximation allowance $\varepsilon_0>0$ before acquisition, sufficient choices are
\[
 n_x\geq\frac{1296}{\varepsilon_0},\qquad
 m\geq\frac{3957}{32\varepsilon_0},\qquad
 n_z\geq\frac{3645}{32\varepsilon_0},\quad 4\mid n_z.
\]
These conservative sufficient resolutions are not a practical recommendation. They make explicit what was absent from a node-count statement alone: fixed original-optimum accuracy can be obtained with resolutions independent of the horizon in these original discounted economies.

Put $S=(n_x+1)^d$, let $A_m=m+1$ in the scalar case and $A_m=(m+1)(m+2)/2$ in the vector case, and let $s\in\{1,2\}$ be the innovation dimension. A dense stored implementation uses $O(TdS^2A_mn_z^s)$ rational arithmetic operations, plus risk sorting and output. Original-owner compilation replaces continuation evaluation by
\[
 O\!\left(T\{dS+SA_mn_z^s\,d[2^d+\log(n_x+1)]\}\right)
\]
operations, plus the same primitive, risk-sorting and output charges. The state and innovation costs remain explicit. Fixed precision is not presumed: a common-denominator realization has rational bit length $O((T+1)\log D)$ for a fixed polynomial-size primitive denominator $D$. Thus operation counts are linear in the horizon at fixed accuracy and dimension, while bit-operation costs also pay for growing exact numbers. This constructive result does not assert that the largest vector run has been independently executed, or that its conservative worst-case resolutions are economically competitive.

\subsection{What is established by each part of the method}\label{sec:comparison95}
Markov risk dynamic programming, conditional-risk duality and policy iteration are established results \citep{ruszczynski2010risk,ruszczynski2014erratum}. Feature-based risk evaluation and structured policy improvement are also established research directions \citep{ruszczynskiyang2025}. We do not claim these objects, or elementary dual centering, as new. Table~\ref{tab:comparison95} separates the prior mathematical tool from the conclusion obtained here.
\begin{table}[htbp]
\caption{Established tools and the additional qualification object}\label{tab:comparison95}
\centering\small
\begin{tabular}{@{}p{.24\linewidth}p{.32\linewidth}p{.37\linewidth}@{}}
\toprule
Tool & Additional information retained & Conclusion in this paper\\\midrule
Risk Bellman recursion & Actual stored continuation and complete feasible controller & State-dependent bounds against the original optimum, not a surrogate objective.\\
Residual comparison & Signed local discrepancies and reference-centered dual penalties & A policy upper inequality and an all-original-action lower inequality; harmless offsets need not be charged twice.\\
Finite occupation duality & Risk-distorted flows and penalties attached to the same continuation & Exact finite comparison with the original $\inf_a\sup_q$ order.\\
Affine support domination & An independently verified error set & Strict safe exclusions; no exclusion at equality and no circular construction of the error set.\\
Convex risk optimization & Primitive curvature, imperfect multiplier residuals and continuous-law bin bounds & Computable whole-action lower witnesses and full-controller upper costs, Theorem~\ref{thm:continuous95}.\\
Feature-based learning & Frozen labels, coefficients, selector and common original-optimum target & A reproducible comparison of representation, policy cost and verification work, not a presumption of neural superiority.\\\bottomrule
\end{tabular}
\end{table}

Three levels must remain distinct. The bounded-function risk theorem is analytical and applies under its stated dual and selection assumptions. The finite affine-risk occupation result has an exact rational implementation. The continuous investment construction has outward numerical enclosures and primitive-specific convexity. Neither of the latter two implementations establishes numerical tractability for all instances of the first. The NBO contribution includes constructive continuations, their attained-policy qualification and their resource account; a validity theorem need not privilege a neural parameterization for that research question to remain meaningful.

\section{A frozen continuation comparison under recursive risk}\label{sec:study95}
We apply the preceding construction to $d=2$ at both declared risk levels. There are four arms: a fitted oblique ReLU continuation, an additive linear-spline continuation, a projected-quadratic continuation, and conventional online evaluation of the original one-date risk Bellman problem. The last arm pays for that calculation; it is not supplied with a free optimal-value oracle. Every fitted label is itself computed from the original terminal primitive and a paid all-action certificate. The terminal continuation is the known $g$, so this last-date fit is consistent with the own-future construction rather than training against an externally supplied optimal continuation.

\subsection{Frozen design and common accuracy}
Each risk level has 1,024 training states, 128 calibration states and 2,048 evaluation states. The three streams are domain-separated before execution. States are drawn from the declared uniform 20-bit dyadic initial-state distribution; the state dynamics and innovations between decisions remain continuous. Each fitted arm uses the same labels, 32 fixed nonlinear features, an unrestricted affine part and nonnegative nonlinear coefficients. Least squares uses the same solver tolerance and iteration cap; readouts are stored on a 30-bit dyadic grid. The source, serialized coefficients and calibration selector are checked before evaluation. No fitting or reselection occurs on the evaluation stream.

The selector minimizes calibration mean cost-interval midpoints with a fixed lexicographic tie rule. It chooses online risk-DP at both risk levels. All twelve pairwise comparisons across the two levels were specified in advance, including those not involving the selected arm. The independent observation is an initial state, not a shock bin or a repeated solver query.

A common own-continuation gap is insufficient for a fair policy comparison. We require both an own-continuation gap at most $1/64$ and an original-optimum complete-policy loss bound at most $1/16$ at every queried initial state. Fitted arms pay for the extra original-continuation lower witness in Theorem~\ref{thm:continuous95}; online risk-DP reuses its already computed witness. Integration may refine from 128 to 4,096 bins, but the proposed first action cannot change. All failures are retained rather than silently replaced by conventional actions. The terminal controller is common to every arm and has the uniform allowance proved above.

\begin{table}[htbp]
\caption{Fitted continuations and whole-square residual enclosures}\label{tab:fits95}
\centering\small
\begin{tabular}{@{}llrrl@{}}
\toprule
$\alpha$ & Continuation & Rank & Training RMSE & $V_1-f$ on the whole square\\\midrule
0.5 & Quadratic & 6 & 0.04122 & [-0.13508, 0.52318]\\
0.5 & ReLU & 35 & 0.01668 & [-0.04373, 0.08665]\\
0.5 & Spline & 35 & 0.03979 & [-0.11690, 0.50298]\\
0.75 & Quadratic & 6 & 0.04166 & [-0.11726, 0.53606]\\
0.75 & ReLU & 35 & 0.01670 & [-0.04599, 0.09043]\\
0.75 & Spline & 35 & 0.03984 & [-0.10431, 0.45189]\\
\bottomrule
\end{tabular}
\par\smallskip\begin{minipage}{.98\linewidth}\footnotesize Each fit uses the same 1,024 labels and 32 nonlinear features plus an affine part. The quadratic feature matrix has rank six, not 35; the experiment matches nominal features and labels, not independent parameter rank. The residual enclosures cover every state in the continuous square.\end{minipage}
\end{table}

The ReLU fits have the smallest training RMSE and the narrowest whole-square residual brackets in both cases. These are distinct observations: training errors alone do not give a whole-domain certificate. The quadratic design has effective rank six although all three fitted arms have 35 stored coefficients; matching feature counts is therefore not matching approximation spaces or statistical degrees of freedom. No rank-matched representation superiority is inferred from this comparison.

\subsection{A common-policy comparison of certificate accounts}
Let $[l,u]$ be the verified whole-square bracket for $V_1-f$, $E=\max\{|l|,|u|\}$, and $\delta(x)$ the own-continuation all-action gap. Two valid original-optimum bounds for the same complete policy are
\[
 B_{\mathrm{norm}}(x)=\delta(x)+\beta(2E+\varepsilon_1),\qquad
 B_{\mathrm{signed}}(x)=\delta(x)+\beta(u-l+\varepsilon_1).
\]
Monotonicity and cash invariance prove these inequalities directly. The joint local-support bound is the paid bound in \eqref{eq:originalgap95}, not a new policy. Table~\ref{tab:accounts95} re-accounts for every existing fitted-arm row under all three bounds. At the fixed whole-square cover, neither scalar account meets $1/16$ on any of these rows, whereas the joint certificate meets it on all 12,288 fitted-arm rows. Thus keeping the economically relevant local continuation information matters for qualification, even though it does not establish a neural policy-cost advantage.

\begin{table}[htbp]
\caption{Three valid certificate accounts for the same frozen fitted policies}\label{tab:accounts95}
\centering\small
\begin{tabular}{@{}llrrrr@{}}\toprule
$\alpha$ & Continuation & Supremum norm & Signed scalar & Joint support & Qualified\\\midrule
0.5 & Quadratic & 0.981430 & 0.617589 & 0.005501 & 2048/2,048\\
0.5 & ReLU & 0.162986 & 0.122750 & 0.021836 & 2048/2,048\\
0.5 & Spline & 0.943607 & 0.581660 & 0.012112 & 2048/2,048\\
0.75 & Quadratic & 1.005610 & 0.612983 & 0.006422 & 2048/2,048\\
0.75 & ReLU & 0.170090 & 0.128429 & 0.022145 & 2048/2,048\\
0.75 & Spline & 0.847773 & 0.521923 & 0.008098 & 2048/2,048\\
\bottomrule\end{tabular}
\par\smallskip\begin{minipage}{.98\linewidth}\footnotesize
Entries are maximum original-optimum policy-loss bounds across the same 2,048 rows. Supremum-norm and signed-scalar bounds use the same frozen whole-square residual cover. Joint support uses the paid original-value witness. No row satisfies $1/16$ with either scalar account; every row does with the joint certificate. This is post-evaluation deterministic re-accounting, not a prospective runtime ablation or an optimality claim for any certificate-construction algorithm.
\end{minipage}\end{table}

This re-accounting was performed after the primary evaluation, with unchanged data, actors and targets. It is a deterministic certificate comparison, not a prospectively registered comparison of optimally tuned centered and uncentered solvers. The scalar accounts pay for the global residual cover; the joint account pays for the original-value witness. Their measured preparation and query work remain recorded separately, so certificate width is not converted into an assumed runtime saving.

\subsection{Original policy loss and economic accuracy}
All 16,384 arm-state records satisfy both qualification targets. The largest original-optimum loss bound among them is less than $0.022145$, below the common $0.0625$ target. This is a statement about every recorded initial state, not an empirical proof of that same tolerance on all possible initial states. The separate residual cover qualifies the whole continuous state square.

\begin{table}[htbp]
\caption{Original nested-risk cost and qualification at 2,048 evaluation states per risk level}\label{tab:cost95}
\centering\small
\begin{tabular}{@{}lllrr@{}}
\toprule
$\alpha$ & Continuation & Mean cost enclosure & Maximum loss bound & Failures\\\midrule
0.5 & Online risk-DP & [0.625634, 0.632502] & 0.003435 & 0\\
0.5 & Quadratic & [0.625754, 0.632622] & 0.005501 & 0\\
0.5 & ReLU & [0.626391, 0.633259] & 0.021836 & 0\\
0.5 & Spline & [0.625792, 0.632660] & 0.012112 & 0\\
0.75 & Online risk-DP & [0.625663, 0.632531] & 0.003435 & 0\\
0.75 & Quadratic & [0.625808, 0.632676] & 0.006422 & 0\\
0.75 & ReLU & [0.626600, 0.633468] & 0.022145 & 0\\
0.75 & Spline & [0.625805, 0.632674] & 0.008098 & 0\\
\bottomrule
\end{tabular}
\par\smallskip\begin{minipage}{.98\linewidth}\footnotesize Cost enclosures here contain numerical uncertainty only. Simultaneous sampling intervals for all paired differences appear in Table~\ref{tab:pairs95}. The maximum loss bound is against the original Bellman optimum; all displayed endpoints are rounded outwards. A failure is a failed own-continuation or original-optimum target, with every row retained.\end{minipage}
\end{table}

The estimand is the initial-distribution mean of each complete policy's original nested-risk value. It is not average realized cash expenditure or a monetary welfare estimate. The fixed margin $1/16$ equals four fifths of the maximal first-date investment charge $h(1/4)=5/64$ in this normalized model. It is a declared numerical economic tolerance; judging it small for a particular application requires an external economic calibration.

All twelve simultaneous empirical-Bernstein intervals lie strictly inside $[-1/16,1/16]$. Thus all four complete policies are equivalent at that prespecified tolerance on the declared initial distribution. None of the intervals excludes zero, so none establishes strict superiority. In particular, the narrower ReLU residual bracket does not translate here into an identified improvement over online risk-DP. The numerical mean-cost brackets in Table~\ref{tab:cost95} are not themselves sampling confidence intervals; the full paired intervals and their sampling components are reported in the appendix.

The numerical component includes original-law integration, terminal observation and optimization, and outward arithmetic. Initial-state discretization is part of the declared evaluation distribution, not an unreported approximation to a continuous uniform initial distribution. There is no transition-law substitution by a state grid, and therefore no hidden state-aggregation allowance in these costs. Sampling and numerical uncertainty are combined as proved in the supplement. The protocol does not permit changing the margin or sample size after examining the evaluation outcomes.

\subsection{Verification work and the value of storage}
\begin{table}[htbp]
\caption{Measured work of the completed evaluation attempt}\label{tab:work95}
\centering\small
\begin{tabular}{@{}llrrr@{}}
\toprule
$\alpha$ & Continuation & Seconds & Terminal proposals & All-action witnesses\\\midrule
0.5 & Online risk-DP & 32.509 & 112,721,920 & 2,048\\
0.5 & Quadratic & 12.254 & 13,631,488 & 4,096\\
0.5 & ReLU & 11.797 & 13,631,488 & 4,096\\
0.5 & Spline & 11.442 & 13,631,488 & 4,096\\
0.75 & Online risk-DP & 30.770 & 112,721,920 & 2,048\\
0.75 & Quadratic & 11.746 & 13,631,488 & 4,096\\
0.75 & ReLU & 11.237 & 13,631,488 & 4,096\\
0.75 & Spline & 11.523 & 13,631,488 & 4,096\\
\bottomrule
\end{tabular}
\par\smallskip\begin{minipage}{.98\linewidth}\footnotesize Every arm processes all 2,048 initial states. Timings include first-action proposals, qualification, the original-value witness, complete-policy cost enclosures and CSV output. They exclude preparation, the independent audit, diagnostic screening and failed earlier execution starts, all reported separately. Counts are instrumented work units, not every arithmetic instruction. This is a single x86-64 execution, not a cross-architecture speed guarantee.\end{minipage}
\end{table}

In the completed execution, the fitted arms require approximately 11.2--12.3 seconds per 2,048-state risk-level workload, compared with approximately 30.8--32.5 seconds for online risk-DP. These measurements include the extra original-optimum witness required of the fitted arms. Preparation for the entire two-risk-level, six-fit comparison takes approximately 8.79 seconds. A separately coded rational audit takes approximately 89.98 seconds and checks the continuous-risk kernels and retained evidence; that verification is a separate paid campaign cost, not zero-cost validation.

These measurements identify a benefit of stored continuations on this workload, not a uniquely neural advantage. Spline and quadratic continuations obtain similar execution reductions, and the fitted arms are not rank matched. The counts identify the repeated terminal computations avoided by storage; they do not identify the cost of every machine instruction. The reported timings are from one x86-64 execution. Two earlier tool-interrupted starts are recorded separately, without recoverable complete cost ledgers. Accordingly the table is not a claim about the aggregate cost of every attempted run, nor a cross-architecture, energy or asymptotic scalability result.

Whole-bin screening excludes between 738 and 746 of 4,096 candidates at $\alpha=1/2$, and between 2,053 and 2,062 at $\alpha=3/4$. After paying for unpruned original-value bounds and the strict test, the screened calculation is slightly slower in every displayed case. The exact unchanged tail values validate the exclusions; they do not establish a computational gain. No-cut rows and all denominators are retained in the appendix. The experiment therefore separates a valid geometric certificate, the amount it removes and the net value of constructing it.

The two-date scalar-control experiment does not replace the inherited two-control, two-innovation horizon frontier. Proposition~\ref{prop:horizon95} now gives explicit fixed-accuracy resolutions and a horizon-uniform construction for both the scalar and original vector economies. The general NBO construction and joint-domain theorem remain available under their hypotheses, while a full independent largest-vector implementation remains a separate numerical obligation. This distinction preserves the general research program without borrowing scalability evidence from a different problem.

\section{Horizon complexity and the price of accuracy}\label{sec:complexity86}
The inherited recursive query experiment exposes descendant growth. The own-future constructive backend has a different dependency graph: once $f_{t+1}$ is stored, a query evaluates that representation and does not solve the remaining Bellman problem again. This distinction changes the horizon account without removing any of the recursive benchmarks.

For each date let $S_t$ be its state-context cover size, $J_{ti}$ its feasible action-net size at node $i$, $M_t$ its certified integration work in primitive representatives, and $E_{t+1}$ the cost of evaluating its stored successor. Construction uses
\begin{equation}\label{eq:dagwork86}
 Q_t=M_t\sum_{i\le S_t}J_{ti},\qquad
 W\le\sum_{t<T}\{Q_t(C_t^{\rm prim}+E_{t+1})+C_t^{\rm compile}\}
       +C_T^{\rm terminal}+C^{\rm output}.
\end{equation}
The formula counts a directed acyclic sequence of stored continuations. A literal descendant solver instead multiplies successor solves by its action and shock branching factors. Neither expression is a lower bound for every possible algorithm. The parameters in \eqref{eq:dagwork86} may themselves grow with horizon if the primitive regularity deteriorates.

\begin{proposition}[Optimal dated error allocation]\label{prop:allocation86}
Suppose a certified constructive implementation has a resource upper account $a_t\eta_t^{-p}$ for attaining total date-$t$ residual allowance $\eta_t>0$, with common $p>0$. Each $a_t$ includes action search, integration, feasibility, compilation and the cost of evaluating the stored successor; it cannot hide another recursive Bellman solve. Put $w_t=B_t>0$ and
\[
 S=\sum_t a_t^{1/(p+1)}w_t^{p/(p+1)}.
\]
Among allocations satisfying $\sum_tw_t\eta_t\le\varepsilon$, the minimum of this resource upper model and its minimizer are
\begin{equation}\label{eq:allocation86}
 W_{\rm model}^*=S^{p+1}\varepsilon^{-p},\qquad
 \eta_t^*=\frac{\varepsilon}{S}(a_t/w_t)^{1/(p+1)}.
\end{equation}
This optimizes the declared upper model, not the unknown fastest algorithm. Rounding each allowance down to a dyadic number within a factor two preserves the error constraint and increases that model by at most $2^p$.

If $a_t\le A$ and $w_t=\beta^t$ for $0<\beta<1$, then
\[
 W_{\rm model}^*\le
 A\varepsilon^{-p}(1-\beta^{p/(p+1)})^{-(p+1)}.
\]
With a maximum permissible allowance $\eta_{\max}$ and positive per-date setup costs, clipping the displayed allocation adds at most
$\sum_t a_t\eta_{\max}^{-p}$ to this upper bound, plus the explicit setup account. For $w_t=1$, the un-clipped bound is $A(T+1)^{p+1}\varepsilon^{-p}$. Thus uniform regularity and reusable stored continuations give polynomial, rather than descendant-tree, horizon dependence at fixed dimensions. They do not give dimension-free computation.
\end{proposition}

The power model follows, for example, when effective state, action and innovation covers and their certified errors have polynomial work bounds, successor evaluation has the separately proved compiled bound, and all moduli are uniformly controlled. The exponent must include all these factors. Tensor compilation in Theorem~\ref{thm:compile49} has a $2^{p_t}$ query factor and cell-location work; a dense neural cone has an $S_{t+1}$ factor. Those are different implementation accounts and must not be interchanged.

\subsection{Operation vectors, arithmetic and reliability}\label{sec:ledger86}
We retain the resource vector
\[
 (N_{\rm primitive},N_{\rm descendant},N_{\rm interval},N_{\rm rational},
 N_{\rm network},B_{\rm read},B_{\rm write},M_{\rm peak}).
\]
Counts are reported before scalar prices are applied. Rational operations additionally record operand bit lengths; an arithmetic-operation count is not a bit-complexity bound. Multiplication and division costs at those lengths enter a bit-cost account. Loading, validation, serialization and shutdown belong to complete process time. Peak memory is a capacity requirement, not an additive flow cost unless its duration or rental convention is specified.

A fitting warning, invalid proposal, zero-progress probe, successful screen, protected fallback and terminal service status are different events. In particular, zero new fits cannot estimate fitting reliability, and a successful fallback cannot be counted as a successful prediction. The retained 336-service recursive catalogue and the nine-controller attained-cost experiment have different units of observation. New mathematical fixtures below establish the operation identities and allocation formula; they are not relabeled as a re-execution of the inherited eight-million-query vector tree or as a hardware calibration. The complete proof of Proposition~\ref{prop:allocation86} is in Supplement~\ref{supp:allocation86}.

\section{Accuracy allocation with endogenous successor cost}\label{sec:coupled89}
The separable allocation in Proposition~\ref{prop:allocation86} applies when its coefficients are fixed independently of the allowances being optimized. A dense successor circuit generally has a size that depends on its own allowance. Charging that size at the preceding date couples the two choices. We now optimize this dependency explicitly.

Let $\eta=(\eta_0,\ldots,\eta_{T-1})$ be positive dated residual allowances. Primitive analysis of a specified stored-continuation implementation supplies the resource upper model
\begin{equation}\label{eq:coupled-work89}
 F(\eta)=F_0+\sum_{j=1}^{J}c_j\prod_{t<T}\eta_t^{-a_{jt}},
 \qquad c_j>0,\quad a_{jt}\ge0.
\end{equation}
The coefficients, exponents and setup term $F_0\ge0$ are fixed before choosing $\eta$. A term $c_t\eta_t^{-p_t}\eta_{t+1}^{-s_{t+1}}$ charges date-$t$ label acquisition and evaluation of its stored successor. Compilation, terminal evaluation and storage-related work may supply other terms. An operation ledger is scalarized only at declared nonnegative resource prices; the model does not estimate monetary prices or energy use.

The feasible accuracy set is
\begin{equation}\label{eq:coupled-budget89}
 \mathcal A=\{\eta:0<\eta_t\le\bar\eta_t,\quad
                    \sum_{t<T}w_t\eta_t\le\varepsilon\},
 \qquad w_t>0,\quad\bar\eta_t>0,\quad\varepsilon>0.
\end{equation}
For the original expected-cost sandwich, $w_t=B_t$. The caps specify where the implementation's error and resource bounds hold. A coordinate with no resource-dependent term can be fixed separately; assume here that every coordinate has a positive exponent in at least one term.

\begin{theorem}[Coupled accuracy optimization with an exact certificate]\label{thm:coupled89}
The upper-model minimization $F_* =\min_{\mathcal A}F(\eta)$ has a unique positive minimizer. It is a convex program in $\eta$ and, after $\eta_t=\exp(x_t)$, also has a convex exponential representation. Standard geometric-programming methods can propose an allocation, but their numerical status is not its certificate.

For any feasible $\eta^0$, write $g=\nabla F(\eta^0)$. Let $\lambda\ge0$ and $\nu_t\ge0$ satisfy
\begin{equation}\label{eq:coupled-certificate89}
 g_t+\lambda w_t+\nu_t\ge0\qquad(t<T).
\end{equation}
Then the independently checkable bounds are
\begin{align}\label{eq:coupled-dual89}
 D(\eta^0,\lambda,\nu)
 &=F(\eta^0)-g\cdot\eta^0-\lambda\varepsilon
                  -\sum_t\nu_t\bar\eta_t,\notag\\
 D(\eta^0,\lambda,\nu)&\le F_*\le F(\eta^0).
\end{align}
At the optimizer, the Karush--Kuhn--Tucker multipliers make these bounds equal. For any candidate, the choice $\nu=0$ and
$\lambda=\max_t\{-g_t/w_t\}$ supplies a valid, possibly loose, lower bound. With integer exponents and rational inputs and candidates, every check is rational arithmetic. For other exponents, certified enclosures of the function and derivatives are required.

Put $P=\max_j\sum_ta_{jt}$. Rounding each allowance down by less than a factor two preserves feasibility and gives
\[
 F(\eta^{\rm dyad})-F_0
 \le2^P\{F(\eta^0)-F_0\}.
\]
The exponent is the total degree of a coupled term, not merely the exponent of the date at which that term is paid.
\end{theorem}

The proof in Supplement~\ref{supp:coupled89} includes attainment, strict convexity and the certificate. Convex duality and geometric programming are established techniques \citep{boyd2004,boyd2007gp}; their role here is to preserve the successor dependency in the NBO's original accuracy account. The certificate bounds the optimum of an explicitly stated resource upper model. It is not a lower bound on the work of every algorithm solving the economic problem.

\begin{corollary}[A constructive horizon bound with coupled evaluation]\label{cor:coupled-horizon89}
Suppose $w_t=1$, $\varepsilon/T\le\min\{1,\min_t\bar\eta_t\}$, $J\le KT$, $c_j\le C$, and the total degrees are bounded by a fixed $P$. Then a feasible stored-continuation implementation satisfies
\[
 F_*\le F_0+KC\,T^{P+1}\varepsilon^{-P}.
\]
For a dense successor account with label exponent $p$ and evaluation exponent $s$, the adjacent-date term has degree $p+s$. Thus a valid polynomial horizon account charges $p+s$, rather than silently holding the successor cost constant. Worsening regularity, growing cover dimension, bit lengths and setup remain explicit inputs; none is bounded by this corollary without its stated hypotheses.
\end{corollary}

\paragraph{A priced two-date allocation.}
Consider the exact upper model
\[
 F(\eta_0,\eta_1)=\eta_0^{-1}
             +(\eta_0\eta_1)^{-1}+2\eta_1^{-2},
 \qquad\eta_0+\eta_1\le\tfrac12,\quad\bar\eta_t=1.
\]
The equal allocation costs $52$. The current numerical appendix gives a rational feasible allocation, its independently recomputed lower certificate and its gap. The example isolates the accuracy dependency algebraically; its resource units are not fitted runtime measurements. It also provides a regression test that rejects a negative multiplier, an infeasible budget or a forged lower bound. In the original economy the same construction applies only after its coefficients have been derived from the actual label, action, integration and stored-circuit accounts.

\paragraph{Policy and resource certificates.}
Qualification determines which policy outputs meet the economic accuracy requirement. The coupled program determines how a specified implementation distributes the corresponding dated allowances. These are different optimizations. A tight resource certificate cannot replace an original-law policy certificate, and a tight policy certificate does not determine hardware time. Keeping both accounts explicit permits constructive improvement without a claim of dimension-free or universally faster computation.

\section{Prospective choice and economically interpretable comparisons}\label{sec:prospective86}
Exact ranking on a completed catalogue and preference before a new workload are different claims. We connect them by a population selection bound and a separately frozen evaluation, retaining the original investment economy and every comparator.

\begin{proposition}[Population regret for a fixed service menu]\label{prop:selection86}
Condition on construction of $m$ complete services, each with original-law cost in $[0,H]$. For $n$ independent new workloads, obtain valid pathwise cost enclosures $[l_{ji},u_{ji}]$. Write $\bar l_j,\bar u_j$ for endpoint means and $q_n=H\sqrt{\log(2m/\delta)/(2n)}$. Select the service minimizing $\lambda\bar u_j+p\cdot w_j+f_j$, with a fixed tie rule. Here $\lambda\ge0$, $p\ge0$, and the nonrandom resource and fee inputs are declared separately. With probability at least $1-\delta$, simultaneously for every such price-fee vector, its expected objective regret is at most
\[
 \lambda\{2q_n+\max_j(\bar u_j-\bar l_j)\}.
\]
A separate evaluation stream with simultaneous pairwise intervals $[l_{jk},u_{jk}]$ gives cost regret at most $\max\{0,\max_k u_{\widehat j,k}\}$ for any calibration-selected service $\widehat j$. This comparison is with the best expected-cost member of the declared menu, not the original Bellman optimum; the original-optimum account remains the separate policy sandwich.
\end{proposition}

\subsection{The frozen independent evaluation}
The prospective protocol preceding the retained evaluation fixes the nine complete two-date investment controllers: the zero controller and both ReLU and Bernstein readouts, under both legacy and calibrated qualification, at acquisition sizes eight and sixteen. Coefficients, complete action tables, constraints and original continuous innovation laws are unchanged. There are 8,192 calibration bins and 131,072 distinct evaluation bins, with forty bits per coordinate. Selection uses the smallest calibration conditioned upper-endpoint mean, with lexical ties, and is fixed before evaluation.

Both independently written whole-bin evaluators run on every policy and bin; rational midpoint fixtures check their hulls. The two estimators are cash cost and analytically conditioned cash cost. Each stream has 36 unordered pairs for each estimator. The empirical-Bernstein error account used in the original study assigns $4\exp(-11)$ to each interval, yielding $288\exp(-11)<0.005$ per stream and $576\exp(-11)<0.01$ jointly. No fitting-success event is inferred. The 139,264 paired workload rows count once, rather than once per policy or per comparison. Reconstructing the frozen streams during the current build adds no independent observations.

\begin{table}[htbp]
\centering
\caption{Independent evaluation: selected fine calibrated Bernstein minus comparator}\label{tab:selection86}
\begin{tabular}{lrr}\toprule
Comparator & Lower endpoint & Upper endpoint\\\midrule
Zero & $-0.177441$ & $-0.170788$\\
Legacy ReLU, size 8 & $-0.010137$ & $-0.005786$\\
Legacy Bernstein, size 8 & $-0.087292$ & $-0.081642$\\
Legacy ReLU, size 16 & $-0.009460$ & $-0.005046$\\
Legacy Bernstein, size 16 & $-0.086152$ & $-0.080530$\\
Calibrated ReLU, size 8 & $-0.009056$ & $-0.004812$\\
Calibrated Bernstein, size 8 & $-0.004823$ & $-0.000716$\\
Calibrated ReLU, size 16 & $-0.003559$ & $0.000573$\\\bottomrule
\end{tabular}
\par\smallskip
{\footnotesize Original-law conditioned cost. The table retains the prespecified R85 evaluation checkpoint; the current reproduction regenerates it from exact integer moments. Endpoints are rounded outwards to six decimal places. Selection uses 8,192 calibration rows, not the 131,072 evaluation rows. The simultaneous evaluation family contains 72 intervals across cash and conditioned estimators. The margin is $1/256$ in normalized cost units.}
\end{table}


Calibration selects the fine calibrated Bernstein controller. In the retained evaluation, its upper expected-cost regret is less than $0.000573$, below the prespecified $1/256$ noninferiority margin. Seven of its eight comparisons have negative upper endpoints. Its comparison with fine calibrated ReLU crosses zero and supports equivalence at that margin, not strict basis superiority. The resolution comparison within calibrated Bernstein is signed. The crossed basis-by-qualification effect is a separate factorial estimand, not a change in the underlying economy. Full endpoint moments, all 144 stream-estimator contrasts, sampling and enclosure components, and all comparisons are retained with the reproduction record.

The margin is a declared acceptable normalized loss. It is not an estimated monetary adoption fee. This evaluation does not replace the historical recursive-query workload, and does not assert that the nine acquired controllers are a prospectively validated neural query proposal on every new regime or horizon. The population construction in Section~\ref{sec:population86} supplies an explicit broader guarantee under its own joint-domain hypotheses; the two empirical objects are not conflated.

\subsection{Uniform price-fee sensitivity}\label{sec:prices86}
Let $\mathcal X$ be a finite set of feasible service assignments, including the original capacity constraints. Let $\mathcal D$ be a nonempty compact polytope containing their joint normalized cost vector. For an assignment $x$, let $w_x$ be its complete resource vector, $\lambda\ge0$ the value of normalized policy loss, $p\ge0$ the resource prices, and $f$ the assignment fee vector. All coordinates refer to explicit units. Put $h_{\mathcal D}(v)=\sup_{c\in\mathcal D}c\cdot v$.

\begin{proposition}[Uniform purchase region]\label{prop:prices86}
Assignment $x$ is optimal for every cost vector in $\mathcal D$ exactly on the region
\begin{equation}\label{eq:price-region86}
 \mathcal R_x=\bigcap_{y\in\mathcal X}
 \{(\lambda,p,f):\lambda h_{\mathcal D}(x-y)
              +p\cdot(w_x-w_y)+f\cdot(x-y)\le0\}.
\end{equation}
This region is polyhedral. Rational support-function upper certificates give an independently checkable sufficient region. Uniformity over a compact polytope of price-fee inputs can be checked at its vertices. The supporting equalities give the breakpoints at which a guaranteed choice can change. A point outside the sufficient region is not automatically evidence that the candidate is inferior.
\end{proposition}

For a single selected service and comparator $k$, a simultaneous cost upper bound $u_{\widehat j,k}$ gives the transparent sufficient condition
\[
 f_{\widehat j}-f_k+p\cdot(w_{\widehat j}-w_k)
 \le-\lambda u_{\widehat j,k}.
\]
The numerical appendix prints every selected-service breakpoint, including negative ones. In particular, the fine-versus-coarse calibrated Bernstein policy comparison supports a positive extra-fee allowance even at zero computation price, whereas a comparison whose upper cost endpoint is positive cannot justify a positive fee solely from policy quality. A resource-driven choice is distinguished from a policy-cost-driven one. Application-specific prices, energy costs and monetary scales require independent calibration; none is manufactured from the normalized tolerance or the favorable timing repetition.

\paragraph{Evidence identities.}
The inherited experiments have distinct observation units. The retained R67 description records 212 service returns and 48 neural fitting warnings; the later training, validation, reuse and vector catalogue records 188 complete processes. The attained-policy factorial study instead has nine controllers and its own independent path rows. The current complete-menu calculation has 2,548 qualification certificates and no new fits. Counts from different catalogues, including the 336-service description in the advisory report, are not added together or assigned to this evaluation without a source-matched manifest. The finer-bin experiment does not erase the older adverse timing and horizon results. Every inherited active section is retained, and the reply identifies the relevant experiment for each contested assertion.

\section{Conclusion}\label{sec:conclusion65}
Neural Bellman Operators connect attained continuation construction to feasible decisions under the original law. Witnesses and residual qualification provide the original-optimum error account; affine calibration identifies a representation-dependent qualification charge. Actionwise cost transfer and exact tolerance frontiers then turn verified optimization accuracy into a sufficient economic stopping condition.

The investment comparison identifies gains from calibrated qualification in both bases and retains the reversal of their earlier ranking. Some cost orderings and equivalence conclusions hold throughout nonzero optimization sublevels without identical actions or new samples. Unresolved coarse classes remain explicit. These conclusions neither assume universal neural superiority nor establish unrestricted fitting reliability.

Conditional purchases and deployment allocations require separately supplied economic scales and charges. Original failed reuse, adverse query timings, vector descendant growth and broader applications retain their records and assumptions. The result is a complete constructive link from an attained approximation to a qualified controller and an economically interpretable accuracy requirement.

\paragraph{Shared coefficients and economic accuracy.}
A qualified output is one controller, not an arbitrary independent choice in every observation cell. Full-menu and shared-readout certificates respect this distinction and supply a constructive refinement of the original actionwise stopping regions. The assignment relaxation need not be an attainable optimizer policy, and every numerical upper bound is therefore identified as an enclosure rather than an observed worst-case loss. The verified common-coefficient restriction strengthens the economic use of the original NBO construction while preserving its original scope and all adverse evidence.

The complete-policy extension propagates changed actions through their own occupation flows. Joint-domain construction establishes uniform shared-operator guarantees under explicit primitive geometry, including endogenous contexts. Stored continuations and optimally allocated dated error budgets give a constructive horizon account distinct from the inherited descendant-tree experiment. Prospective selection and uniform price-fee regions then separate statistical policy quality from the cost of computation. Nonlinear preferences and strategic deviations retain their own valid comparison principles, rather than inheriting an unsupported interpretation of an expected-cost table.


\paragraph{The original law remains the benchmark.}
Pattern-consistent occupations and transport-qualified residuals provide complementary complete-policy guarantees. The former retains shared decisions across the contingent plan; the latter prices the action-value effect of numerical transition approximation, including singular laws. Both extend the original NBO construction and preserve its original-optimum interpretation. Their proofs identify every required regularity, qualification and arithmetic condition. The prospective evidence and coupled accuracy account connect those conditions to measured or explicitly priced resources without replacing unmeasured costs by assumed savings.


Signed qualification further separates a harmless global offset from a state-dependent approximation that changes complete-policy incentives. The recursive order proof retains the original preference primitive, and the explicit tail-risk and entropic moduli price transition information at the stated economic accuracy. Its sharpness examples and exhaustive finite-policy diagnostics support verification without being counted as new economic observations. This extends the original NBO argument while leaving empirical algorithm comparisons, application-specific prices and the largest vector reimplementation as distinct evidence requirements.


The centered-risk extension retains where an error occurs and which risk kernels can make it consequential. Its exact finite occupation representation and safe kernel exclusions sharpen the same original-optimum account. Full construction and arithmetic costs remain separate from the size of a verified certificate; the unchanged empirical comparisons are not converted into neural-superiority claims by this analytical improvement.

The continuous-law construction makes the required lower and upper inequalities executable in the investment model. Its primitive proof is distinct from the numerical fit, and its common original-optimum target is distinct from agreement with a model's own continuation. The resulting comparison identifies where a stored representation improves approximation and where the full policy or its computation does not improve. This is the evidential distinction required to develop Neural Bellman Operators as a numerical economic method.

\bibliographystyle{ecta-fullname}
\bibliography{references92-canonical}
\end{document}
